\subsection{整闭整环的分解与惯性}

引理 4. 设 A 为整闭整环，K 为其分式域，p 为 K 的特征指数，K' 为 K 的一个纯不可分扩张，A' 为 K' 中包含 A 且在 A 上整的子环。对 A 的每个素理想 p，存在唯一的素理想 \(\mathfrak{p}'\)，即 \(\mathbf{A}'\) 中位于 \(\mathfrak{p}\) 之上的素理想，且 \(\mathfrak{p}'\) 是由满足如下条件的 \(x \in \mathbf{A}'\) 组成的集合：存在整数 \(m \geqslant 0\) 使得 \(x^{p^m} \in \mathfrak{p}\)。

\(\mathfrak{p}'\) 的存在性由第 1 节定理 1 得出。若 \(x \in \mathfrak{p}'\)，则存在 \(m \geqslant 0\) 使 \(x^{p^m} \in \mathbf{K}\)，于是 \(x^{p^m} \in \mathbf{A}\)（因 \(\mathbf{A}\) 整闭），从而 \(x^{p^m} \in \mathfrak{p}' \cap \mathbf{A} = \mathfrak{p}\)。反之，若 \(x \in \mathbf{A}'\) 满足 \(x^{p^m} \in \mathfrak{p} \subset \mathfrak{p}'\)，则 \(x \in \mathfrak{p}\)（因 \(\mathfrak{p}'\) 素）。

注 (1). 由 § 1 第 3 节命题 11 的推论可知，A 在 \(\mathbf{K}'\) 中的整闭包是由满足如下条件的 \(x \in \mathbf{K}'\) 组成的集合：存在 \(m \geqslant 0\) 使 \(x^{p^m} \in \mathbf{A}\)（《代数》第 V 章 § 8 第 1 节命题 1）。

命题 6. 设 A 为整闭整环，K 为其分式域，K' 为 K 的拟 Galois 扩张，A' 为 A 在 K' 中的整闭包。那么：

\begin{enumerate}
\def\labelenumi{(\roman{enumi})}
\item
  对 A 的每个素理想 p，K' 的 K-自同构群 G 传递地作用于 A' 的位于 p 之上的素理想所成的集合。
\item
  对 \(\mathfrak{p}'\)（\(\mathbf{A}'\) 的任意素理想），分式域 \(k'\)（即 \(\mathbf{A}' / \mathfrak{p}'\) 的分式域）是 \(k\)（\(\mathbf{A} / (\mathbf{A} \cap \mathfrak{p}')\) 的分式域）的拟 Galois 扩张，且典范同态 \(\sigma \mapsto \bar{\sigma}\)（从 \(\mathcal{G}^{\mathbb{Z}}(\mathfrak{p}')\) 到群 \(\Gamma\)，后者由 \(k\)-自同构（\(k'\) 上的）组成）经取商给出 \(\mathcal{G}^{\mathbb{Z}}(\mathfrak{p}') / \mathcal{G}^{\mathbb{T}}(\mathfrak{p}')\) 到 \(\Gamma\) 上的一个双射。
\end{enumerate}

\begin{enumerate}
\def\labelenumi{(\Alph{enumi})}
\item
  先设 \(\mathbf{K}'\) 是 \(\mathbf{K}\) 的有限 Galois 扩张。此时 \(\mathbf{A} = \mathbf{A}' \cap \mathbf{K}\)（因 \(\mathbf{A}\) 整闭），故 \(\mathbf{A}\) 是 \(\mathcal{G}\) 在 \(\mathbf{A}'\) 中的不变量环。由于 \(\mathcal{G}\) 有限，命题在这一情形由第 2 节定理 2 得出。
\item
  其次设 \(\mathbf{K}'\) 是 \(\mathbf{K}\) 的任意 Galois 扩张。此时 \(\mathbf{K}'\) 是一个右定向族 \((\mathrm{K}_{\alpha})_{\alpha \in I}\)（由 \(\mathbf{K}\) 的有限 Galois 扩张组成）的并。为证 (i)，考虑 \(\mathfrak{p}'\)、\(\mathfrak{q}'\)，即 \(\mathbf{A}'\) 的位于 \(\mathfrak{p}\) 之上的两个素理想。对每个 \(\alpha \in I\)，\(\mathfrak{p}' \cap \mathrm{K}_{\alpha}\) 与 \(\mathfrak{q}' \cap \mathrm{K}_{\alpha}\) 是 \(\mathbf{A}' \cap \mathrm{K}_{\alpha}\) 的位于 \(\mathfrak{p}\) 之上的两个素理想。由于 \(\mathbf{A}' \cap \mathrm{K}_{\alpha}\) 是 \(A\) 在 \(\mathrm{K}_{\alpha}\) 中的整闭包，且在 \(\mathrm{K}_{\alpha}\) 上 \(\mathcal{G}\) 的元素的限制构成 \(\mathbf{K}\) 的 K-自同构群，故由情形 (A) 存在 \(\sigma \in \mathcal{G}\) 使 \(\sigma. (\mathfrak{p}' \cap \mathrm{K}_{\alpha}) = \mathfrak{q}' \cap \mathrm{K}_{\alpha}\)。设 \(\mathcal{E}_{\alpha}\) 为具有上述性质的 \(\sigma \in \mathcal{G}\) 组成的集。取 \(\sigma \in \mathcal{G} - \mathcal{E}_{\alpha}\)；则对每个 \(\tau \in \mathcal{G}\)，它使 \(\mathrm{K}_{\alpha}, (\sigma \tau). (\mathfrak{p}' \cap \mathrm{K}_{\alpha}) = \sigma. (\mathfrak{p}' \cap \mathrm{K}_{\alpha}) \neq \mathfrak{q}' \cap \mathrm{K}_{\alpha}\) 成立，因而 \(\sigma \tau \in \mathcal{G} - \mathcal{E}_{\alpha}\)。由此可知 \(\mathcal{E}_{\alpha}\) 在拓扑 Galois 群 \(\mathcal{G}\) 中是闭的（《代数》第 V 章附录 II 第 1 节），且族 \((\mathcal{E}_{\alpha})_{\alpha \in I}\) 显然是左定向的。由于 \(\mathcal{G}\) 紧（同前引，第 2 节，命题 3）且诸 \(\mathcal{E}_{\alpha}\) 非空，交 \(\mathcal{E}\)（即族 \((\mathcal{E}_{\alpha})\) 的交）非空，且有 \(\sigma. \mathfrak{p}' = \mathfrak{q}'\)（对每个 \(\sigma \in \mathcal{E}\)），由此得 (i)。
\end{enumerate}

为证 (ii)，注意 \(k'\) 是右定向族 \((k_{\alpha})_{\alpha \in \mathbf{I}}\) 的并，其中 \(k_{\alpha}\) 是 \((\mathrm{A}' \cap \mathrm{K}_{\alpha}) / (\mathfrak{p}' \cap \mathrm{K}_{\alpha})\) 的分式域。由于每个 \(k_{\alpha}\) 由 (A) 是 \(k\) 的拟 Galois 扩张，故 \(k'\) 亦然（《代数》第 V 章 § 6 第 3 节命题 8）。另一方面，设 \(u\) 是一个 \(k\)-自同构，作用于 \(k'\)，并设 \(\pi': \mathrm{A}' \to \mathrm{A}' / \mathfrak{p}'\) 为典范同态。对 \(\mathrm{A}' \cap \mathrm{K}_{\alpha}\) 应用第 2 节定理 2 可知，对每个 \(\alpha\) 存在一个非空集 \(\mathcal{F}_{\alpha}\)，由具有下述性质的元素 \(\sigma \in \mathcal{G}\) 组成：\(\sigma \cdot (\mathfrak{p}' \cap \mathrm{K}_{\alpha}) = \mathfrak{p}' \cap \mathrm{K}_{\alpha}\)，且 \(u(\pi'(x)) = \pi'(\sigma.x)\) 对一切 \(x \in \mathrm{A}' \cap \mathrm{K}_{\alpha}\) 成立。与上面一样可以看出 \(\mathcal{F}_{\alpha}\) 在 \(\mathcal{G}\) 中闭，且由于 \((\mathcal{F}_{\alpha})\) 是左定向集，其交 \(\mathcal{F}\) 非空。显然对 \(\sigma \in \mathcal{F}\)，有 \(\sigma \in \mathcal{G}^{\mathbb{Z}}(\mathfrak{p}')\) 且 \(\bar{\sigma} = u\)，这就在此情形完成了 (ii) 的证明。

\begin{enumerate}
\def\labelenumi{(\Alph{enumi})}
\setcounter{enumi}{2}
\tightlist
\item
  一般情形。G 的不变量域 \(K_{1}\) 是 K 的纯不可分扩张（《代数》第 V 章 § 10 第 9 节命题 14）；因此存在位于 p 之上的唯一素理想 \(p_{1}\)，即 \(A_{1} = A' \cap K_{1}\) 的素理想（引理 4）。若 \(p'\) 与 \(q'\) 是 \(A'\) 的位于 p 之上的两个素理想，则它们都位于 \(p_{1}\) 之上；由于 \(K'\) 是 \(K_{1}\) 的 Galois 扩张且 \(A' \cap K_{1}\) 整闭（§ 1 第 2 节命题 7 与命题 8 的推论），由 (B) 存在 \(\sigma \in G\) 使 \(\sigma \cdot p' = q'\)；由此得 (i)。另一方面，显然 \(k_{1}\)（即 \(A_{1}/p_{1}\) 的分式域）是 k 的纯不可分扩张（因 A 整闭）；由于 \(k'\) 由 (B) 是 \(k_{1}\) 的拟 Galois 扩张，故 \(k'\) 是 k 的拟 Galois 扩张，因为 \(k'\) 到 \(k'\) 的一个代数闭扩张的每个 k-同构都是 \(k_{1}\)-同构。最后这个注记结合 (B) 还表明，\(k'\) 的每个 k-自同构都具有形状 \(\bar{\sigma}\)，其中 \(\sigma \in \mathcal{G}^{\mathbb{Z}}(\mathfrak{p}')\)，这就完成了 (ii) 的证明。
\end{enumerate}

注

\begin{enumerate}
\def\labelenumi{(\arabic{enumi})}
\setcounter{enumi}{1}
\tightlist
\item
  设 \(\mathbf{K}'\) 是 \(\mathbf{K}\) 的 Galois 扩张，并保留命题 6 的证明中的记号；对每个 \(\alpha\)，设 \(\mathcal{G}_{\alpha}^{\mathbb{Z}}\)（相应地 \(\mathcal{G}_{\alpha}^{\mathrm{T}}\)）为 \(\mathcal{G}\) 的由诸 \(\sigma\)（其在 \(\mathrm{A}' \bigcap \mathrm{K}_{\alpha}\) 上的限制属于 \(\mathcal{G}^{\mathbb{Z}}(\mathfrak{p}' \bigcap \mathrm{K}_{\alpha})\)（相应地 \(\mathcal{G}^{\mathrm{T}}(\mathfrak{p}' \bigcap \mathrm{K}_{\alpha})\)）。命题 6 的证明表明这些子群在 \(\mathcal{G}\) 中闭，且
\end{enumerate}

\[
\mathcal {G} ^ {\mathrm{Z}} \left(\mathfrak {p} ^ {\prime}\right) = \bigcap_ {\alpha} \mathcal {G} _ {\alpha} ^ {\mathrm{Z}} \quad \text { and } \quad \mathcal {G} ^ {\mathrm{T}} \left(\mathfrak {p} ^ {\prime}\right) = \bigcap_ {\alpha} \mathcal {G} _ {\alpha} ^ {\mathrm{T}}.
\]

而且，在 \(A^{\prime} \cap K_{\alpha}\) 上，G（相应地 \(G_{\alpha}^{T}\)）的元素的限制所成的集就是整个群 \(\mathcal{G}^{\mathtt{Z}}(\mathfrak{p}^{\prime} \cap \mathtt{K}_{\alpha})\)（相应地 \(\mathcal{G}^{\mathtt{T}}(\mathfrak{p}^{\prime} \cap \mathtt{K}_{\alpha})\)），因为 \(K_{\alpha}\) 的每个 K-自同构都可延拓为 G 的一个元素。

在同一假设下，环 \(\mathbf{A}^{\mathbf{Z}}(\mathfrak{p}')\)（相应地 \(\mathbf{A}^{\mathbf{T}}(\mathfrak{p}')\)）是定向族 \(\mathbf{A}^{\mathbf{Z}}(\mathfrak{p}' \cap \mathbf{K}_{\alpha})\)（相应地 \(\mathbf{A}^{\mathbf{T}}(\mathfrak{p}' \cap \mathbf{K}_{\alpha})\)）的并：事实上，每个 \(x \in \mathbf{A}^{\mathbf{Z}}(\mathfrak{p}')\)（相应地每个 \(x \in \mathbf{A}^{\mathbf{T}}(\mathfrak{p}')\)）属于某个 \(\mathbf{K}_{\alpha}\)，且由前所述存在 \(\beta\) 使 \(\mathbf{K}_{\alpha} \subset \mathbf{K}_{\beta}\)，并且到 \(\mathbf{A}' \cap \mathbf{K}_{\alpha}\) 上，\(\mathcal{G}^{\mathbf{Z}}(\mathfrak{p}')\)（相应地 \(\mathcal{G}^{\mathbf{T}}(\mathfrak{p}')\)）的元素的限制与到 \(\mathbf{A}' \cap \mathbf{K}_{\alpha}\) 上 \(\mathcal{G}^{\mathbf{Z}}(\mathfrak{p}' \cap \mathbf{K}_{\beta})\)（相应地 \(\mathcal{G}^{\mathbf{T}}(\mathfrak{p}' \cap \mathbf{K}_{\beta})\)）的元素的限制相同，群 \(\mathcal{G}^{\mathbf{Z}}(\mathfrak{p}' \cap \mathbf{K}_{\alpha})\) 与 \(\mathcal{G}^{\mathbf{T}}(\mathfrak{p}' \cap \mathbf{K}_{\beta})\) 是有限的；因此 \(x\) 属于 \(\mathbf{A}^{\mathbf{Z}}(\mathfrak{p}' \cap \mathbf{K}_{\beta})\)（相应地 \(\mathbf{A}^{\mathbf{T}}(\mathfrak{p}' \cap \mathbf{K}_{\beta})\)）。

推论 1. 在命题 6 的假设下，设 \(f\) 是从 \(\mathbf{A}\) 到一个域 \(\mathbf{L}\) 的同态，\(g_1\)、\(g_2\) 是从 \(\mathbf{A}'\) 到 \(\mathbf{L}\) 且延拓 \(f\) 的两个同态。那么存在一个 \(\mathbf{K}\)-自同构 \(\sigma\)，即 \(\mathbf{K}'\) 的自同构，使 \(g_1 = g_2 \circ \sigma\)。

从命题 6 出发给出本推论的证明，与从定理 2 出发证明该定理的推论时的做法相同。

推论 2. 设 A 为整闭整环，K 为其分式域，K' 为 K 的有限代数扩张，A' 为 K' 中包含 A 且在 A 上整的子环。

\begin{enumerate}
\def\labelenumi{(\roman{enumi})}
\item
  对 \(\mathfrak{p}\)（\(\mathbf{A}\) 的每个素理想），\(\mathbf{A}'\) 的位于 \(\mathfrak{p}\) 之上的素理想所成的集是有限的。\\
\item
  若 \(\mathfrak{p}'\) 是 \(\mathbf{A}'\) 的位于 \(\mathfrak{p}\) 之上的一个素理想，则 \(\mathbf{A}' / \mathfrak{p}'\) 的每个元素都具有次数 \(\leqslant [\mathbf{K}' : \mathbf{K}]\)，这里是对 \(\mathbf{A} / \mathfrak{p}\) 的分式域而言。
\item
  设 \(\mathbf{K}''\) 为 \(\mathbf{K}\) 的拟 Galois 扩张，它是在 \(\mathbf{K}'\) 的一个代数闭包中由 \(\mathbf{K}'\) 生成的；设 \(\mathbf{A}''\) 为 \(\mathbf{A}\) 在 \(\mathbf{K}''\) 中的整闭包。域 \(\mathbf{K}''\) 是 \(\mathbf{K}\) 的有限扩张（《代数》第 V 章 § 6 第 3 节命题 9 的推论 1），因而其 K-自同构群有限；由此可知 \(\mathbf{A}''\) 的位于 \(p\) 之上的素理想所成的集是有限的（命题 6 (i)）。另一方面，由于 \(\mathbf{A}''\) 在 \(\mathbf{A}\) 上整，映射 \(p'' \mapsto p'' \cap \mathbf{A}'\)（从 \(\mathbf{A}''\) 的位于 \(p\) 之上的素理想之集到 \(\mathbf{A}'\) 的位于 \(x\) 之上的素理想之集）是满射（第 1 节定理 1）。
\item
  任意元素 \(x' \in A'\) 的极小多项式（在 K 上的）的系数都属于 A（§ 1 第 3 节命题 10 的推论）；对这个多项式的系数施加典范同态 \(\pi': A' \to A'/p'\)，便得到一个系数在 A/p 中、次数 \(\leqslant [K': K]\) 的整依赖方程，它为模 \(p'\) 之下的 \(x\) 的剩余类所满足；由此即得结论。
\end{enumerate}

推论 3. 在推论 2 的假设与记号下，若 A 是半局部环，则 A' 亦然。

对每个极大理想 \(m'\)（属于 \(A'\)），\(m' \cap A\) 都是 A 的一个极大理想（第 1 节命题 1）；由于假设 A 的极大理想所成的集有限，本推论随即由推论 2 得出。

推论 4. 设 A 为整闭整环，K 为其分式域，K' 为 K 的 Galois 扩张，A' 为 A 在 K' 中的整闭包，p' 为 A' 的一个素理想，p = A ∩ p'，且 k 与 k' 分别为 A/p 与 A'/p' 的分式域。那么：

\begin{enumerate}
\def\labelenumi{(\roman{enumi})}
\item
  \(\mathbf{A}^{\mathbf{z}} / (\mathfrak{p}'\cap \mathbf{A}^{\mathbf{z}})\) 的分式域等于 \(k\)，而 \(\mathbf{A}^{\mathbf{z}}\) 相对于 \(\mathfrak{p}'\cap \mathbf{A}^{\mathbf{z}}\) 的局部环的极大理想由 \(\mathfrak{p}\) 生成。
\item
  分式域 \(k^{\mathbf{T}}\)（即 \(\mathbf{A}^{\mathbf{T}} / (\mathfrak{p}' \cap \mathbf{A}^{\mathbf{T}})\) 的分式域）是 \(k\) 的包含在 \(k'\) 中的最大可分扩张。
\end{enumerate}

环 A 是 K' 在 K 上的 Galois 群在 A' 中的不变量环；若 K' 在 K 上是有限次的，则本推论由第 2 节的命题 4 与命题 5 得出。现在考虑一般情形，此时 K' 是由 K 的有限 Galois 扩张组成的右定向集 (K \(_{\alpha}\) ) 的并。那么：

\begin{enumerate}
\def\labelenumi{(\roman{enumi})}
\tightlist
\item
  若 \(x, y\) 是 \(\mathbf{A}^z\) 的两个元素，其中 \(y \notin \mathfrak{p}'\)，则存在指标 \(\alpha\) 使 \(x\) 与 \(y\) 属于 \(\mathbf{A}^z(\mathfrak{p}' \cap \mathbf{K}_\alpha)\)（注 2）；由第 2 节的命题 4，存在 \(x_0, y_0\)（在 \(\mathbf{A}\) 中），满足 \(y_0 \notin \mathfrak{p}'\) 且 \(xy_0 - x_0y \in \mathfrak{p}'\)，这就证明了 (i) 的第一个断言；若进一步 \(x \in \mathfrak{p}'\)，则可假定 \(y_0\) 满足
\end{enumerate}

\[
x y _ {0} \in \mathfrak {p A} ^ {\mathrm{z}} (\mathfrak {p} ^ {\prime} \cap \mathrm{K} _ {\alpha}) \subset \mathfrak {p A} ^ {\mathrm{z}} (\mathfrak {p} ^ {\prime}),
\]

这就证明了 (i) 的第二个断言。

\begin{enumerate}
\def\labelenumi{(\roman{enumi})}
\setcounter{enumi}{1}
\tightlist
\item
  现在设 \(x \in \mathbf{A}^{\mathrm{T}}\)；存在 \(\alpha\) 使 \(x \in \mathbf{A}^{\mathrm{T}}(\mathfrak{p}' \cap \mathbf{K}_{\alpha})\)（注 2），且第 2 节的命题 5 表明，剩余类 \(\bar{x}\)（即 \(x \bmod. (\mathfrak{p}' \cap \mathbf{K}_{\alpha} \cap \mathbf{A}^{\mathrm{T}})\) 的类）在 \(k\) 上是代数的且可分的；从而模 \((\mathfrak{p}' \cap \mathbf{A}^{\mathrm{T}})\) 之下 \(x\) 的剩余类更有理由在 \(k\) 上可分；为完成本推论的证明，只需证明 \(k'\) 是 \(k^{\mathrm{T}}\) 的纯不可分扩张。而 \(k'\) 是由分式域 \(k_{\alpha}\)（即环 \((\mathbf{A}' \cap \mathbf{K}_{\alpha}) / (\mathfrak{p}' \cap \mathbf{K}_{\alpha})\) 的分式域）组成的右定向族的并。于是由命题 5 可得：若 \(k'\) 的一个元素属于 \(k_{\alpha}\)，则它在如下商环的分式域上是纯不可分的
\end{enumerate}

\[
\mathbf {A} ^ {\mathrm{T}} \left(\mathfrak {p} ^ {\prime} \cap \mathbf {K} _ {\alpha}\right) / \left(\mathfrak {p} ^ {\prime} \cap \mathbf {A} ^ {\mathrm{T}} \left(\mathfrak {p} ^ {\prime} \cap \mathbf {K} _ {\alpha}\right)\right)
\]

从而更有理由在 \(k^{\mathbf{T}}\) 上是纯不可分的（依据注 2）。

定义 5. 在命题 6 的假设与记号下，不变量域 \(\mathbf{K}^{\mathrm{Z}}(\mathfrak{p}')\)（相应地 \(\mathbf{K}^{\mathrm{T}}(\mathfrak{p}')\)，即群 \(\mathcal{G}^{\mathrm{Z}}(\mathfrak{p}')\)（相应地 \(\mathcal{G}^{\mathrm{T}}(\mathfrak{p}')\)）在域 \(\mathbf{K}'\) 中的不变量域）称为 \(\mathfrak{p}'\) 关于 \(\mathbf{K}\) 的分解域（相应地惯性域）。

我们也用 \(K^{z}\)（相应地 \(K^{T}\)）代替 \(\mathbf{K}^{\mathbf{Z}}(\mathfrak{p}')\)（相应地 \(\mathbf{K}^{\mathrm{T}}(\mathfrak{p}')\)）来记。由 §1 第 9 节命题 23 可知，\(K^{z}\)（相应地 \(K^{T}\)）是环 \(A^{z}\)（相应地 \(A^{T}\)）的分式域；\(A^{z}\)（相应地 \(A^{T}\)）是 A 在 \(K^{z}\)（相应地 \(K^{T}\)）中的整闭包。

\textbf{注}\par

\begin{enumerate}
\def\labelenumi{(\arabic{enumi})}
\setcounter{enumi}{2}
\tightlist
\item
  在命题 6 的推论 4 的条件下，并设 \([\mathbf{K}':\mathbf{K}]\) 有限，位于 \(p\) 之上的不同素理想的个数是 \([\mathbf{K}^{\mathbf{Z}}:\mathbf{K}]\)，由 Galois 理论，此次数等于指数 \((\mathcal{G}:\mathcal{G}^{\mathbf{Z}})\)；此外，由 Galois 理论可得
\end{enumerate}

\[
[ \mathbf {K} ^ {\mathrm{T}}: \mathbf {K} ^ {\mathrm{Z}} ] = (\mathcal {G} ^ {\mathrm{Z}}: \mathcal {G} ^ {\mathrm{T}}) = [ k ^ {\mathrm{T}}: k ].\tag{6}
\]

\begin{enumerate}
\def\labelenumi{(\arabic{enumi})}
\setcounter{enumi}{3}
\tightlist
\item
  设 A 为整闭整环，K 为其分式域，\(K'\) 为 K 的有限代数扩张，\(A'\) 为 A 在 \(K'\) 中的整闭包。那么，对 A 的每个素理想 p，\(A'\) 的位于 p 之上的素理想的个数至多为 \([K':K]_{s}\)（即 \(K'\) 在 K 上的次数的可分部分）。我们可以先把注意力限制到 \(K'\) 是 K 的可分扩张的情形，因为一般地 \(K'\) 是 K 的最大可分扩张 \(K_{0}\)（即包含在 \(K'\) 中的那个）的纯不可分扩张，按定义 \([K':K]_{s} = [K_{0}:K]\)，而且若 \(A_{0}\) 是 A 在 \(K_{0}\) 中的整闭包，则 \(A_{0}\) 与 \(A'\) 的素理想一一对应（引理 4）。因此设 \(K'\) 在 K 上可分，并设 N 为在 K 的一个代数闭包中由 \(K'\) 生成的 K 的 Galois 扩张，G 为其 Galois 群，B 为 A 在 N 中的整闭包，P 为 B 的位于 p 之上的一个素理想。设 H 为 N 在 \(K'\) 上的 Galois 群，\(G^{z}\) 为 P 的分解群；B 的位于 p 之上的素理想就是诸 s.P，其中 \(s \in G\)（第 2 节定理 2），而关系 s.P = s'.P 意味着 \(s' = sg\)，其中 \(g \in G^{z}\)。另一方面，要使 s.P ∩ \(K' = s'\).P ∩ \(K'\)，必须且只需 \(s'\).P = ts.P，其中 \(t \in H\)（第 2 节定理 2），最终得 \(s' = tsg\)，其中 \(t \in H\) 且 \(g \in G^{z}\)。于是 A' 的位于 p 之上的素理想的个数等于 G 按等价关系「存在 \(t \in H\) 与 \(g \in G^{z}\) 使 \(s' = tsg\)」（在 s 与 \(s'\) 之间）所分的类的个数；显然这个数至多等于指数 \((\mathcal{G}:\mathcal{H})\)，即 H 在 G 中的右陪集数，而由 Galois 理论 \((\mathcal{G}:\mathcal{H}) = [\mathrm{K}':\mathrm{K}]\)。
\end{enumerate}

命题 7. 设 A 为整闭整环，K 为其分式域，K' 为 K 的 Galois 扩张，G 为其 Galois 群，A' 为 A 在 K' 中的整闭包，p' 为 A' 的素理想，且 p = A ∩ p'。最后，设 L 为 K' 中包含 K 的子域，并令 p(L) = p' ∩ L。

\begin{enumerate}
\def\labelenumi{(\roman{enumi})}
\item
  \(\mathfrak{p}'\) 关于 \(\mathbf{L}\) 的分解域（相应地惯性域）是 \(\mathbf{L}(\mathbf{K}^{\mathbf{Z}})\)（相应地 \(\mathbf{L}(\mathbf{K}^{\mathbf{T}})\)）；若进一步 \(\mathbf{L}\) 是 \(\mathbf{K}\) 的 Galois 扩张，则 \(\mathfrak{p}(\mathbf{L})\) 关于 \(\mathbf{K}\) 的分解域是 \(\mathbf{L} \cap \mathbf{K}^{\mathbf{Z}}\)。
\item
  若 \(\mathbf{L}\) 包含于 \(\mathbf{K}^{\mathbf{Z}}\)，则 \(\mathbf{A} / \mathfrak{p}\) 与 \((\mathbf{A}' \cap \mathbf{L}) / \mathfrak{p}(\mathbf{L})\) 有相同的分式域，并且在局部环 \(\mathbf{A}' \cap \mathbf{L}\)（对应于素理想 \(\mathfrak{p}(\mathbf{L})\)）中，极大理想由 \(\mathfrak{p}\) 生成。反之，若这两个条件成立，且进一步 \(\bigcap_{n \geqslant 0} \mathfrak{p}^{n} \mathbf{A}_{\mathfrak{p}'}' = 0\)，则 \(\mathbf{L}\) 包含于 \(\mathbf{K}^{\mathbf{Z}}\)。
\item
  若 \(\mathcal{H}\) 是 \(\mathcal{G}\) 的使 \(\mathbf{L}\) 的元素保持不变的那个子群，则显然 \(\mathfrak{p}'\) 关于 \(\mathbf{L}\) 的分解群（相应地惯性群）是 \(\mathcal{G}^{\mathrm{Z}}\cap \mathcal{H}\)（相应地 \(\mathcal{G}^{\mathrm{T}}\cap \mathcal{H}\)），而当 \(\mathbf{K}'\) 是 \(\mathbf{K}\) 的有限 Galois 扩张时，第一个断言由 Galois 理论得出（《代数》第 V 章 § 10 第 6 节定理 3 的推论 1）；一般情形则由如下事实得出：在注 2 的记号下，\(\mathbf{A}^{\mathrm{Z}}\)（相应地 \(\mathbf{A}^{\mathrm{T}}\)）是诸 \(\mathbf{A}^{\mathrm{Z}}(\mathfrak{p}'\cap \mathbf{K}_{\alpha})\)（相应地 \(\mathbf{A}^{\mathrm{T}}(\mathfrak{p}'\cap \mathbf{K}_{\alpha})\)）的并：每个元素 \(x\in \mathbf{K}'\) 属于某个 \(\mathbf{K}_{\alpha}\)，而若它在 \(\mathcal{G}^{\mathrm{Z}}(\mathfrak{p}')\cap \mathcal{H}\)（相应地 \(\mathcal{G}^{\mathrm{T}}(\mathfrak{p}')\cap \mathcal{H}\)）下不变，则它也在 \(\mathcal{G}^{\mathrm{Z}}(\mathfrak{p}'\cap \mathbf{K}_{\beta})\cap \mathcal{H}\)（相应地 \(\mathcal{G}^{\mathrm{T}}(\mathfrak{p}'\cap \mathbf{K}_{\beta})\cap \mathcal{H}\)）下不变，其中 \(\beta\) 是某个适当的指标；于是由论证的开头部分，它属于
\end{enumerate}

\[
\mathbf {L} \left(\mathrm{K} ^ {\mathrm{Z}} \left(\mathfrak {p} ^ {\prime} \cap \mathrm{K} _ {\alpha}\right)\right) \subset \mathbf {L} \left(\mathrm{K} ^ {\mathrm{Z}}\right) (\text { resp. } \mathbf {L} \left(\mathrm{K} ^ {\mathrm{T}} \left(\mathfrak {p} ^ {\prime} \cap \mathrm{K} _ {\beta}\right)\right) \subset \mathbf {L} \left(\mathrm{K} ^ {\mathrm{T}}\right)).
\]

现在设 L 是 K 的 Galois 扩张；于是每个 \(\sigma\in G^{z}\) 在 L 上的限制都保持 \(\mathfrak{p}(L)=\mathfrak{p}^{\prime}\cap L\) 不变，从而属于 \(\mathfrak{p}(L)\) 关于 K 的分解群。反之，设 \(\tau\) 是属于该群的 L 的一个自同构，并设 \(\sigma\) 是 \(\tau\) 到 \(K^{\prime}\) 的一个 K-自同构上的延拓；记 \(q^{\prime}=\sigma.p^{\prime}\)。由于 \(p^{\prime}\) 与 \(q^{\prime}\) 都位于 \(\mathfrak{p}(L)\) 之上，故存在自同构 \(\rho\in H\) 使 \(q^{\prime}=\rho.p^{\prime}\)，从而 \(\rho^{-1}\sigma\in G^{z}\)，且 \(\tau\) 是 \(\rho^{-1}\sigma\) 在 L 上的限制；换句话说，\(\mathfrak{p}(L)\) 关于 K 的分解群恰为诸自同构 \(\sigma\in G^{z}\) 在 L 上的限制所成的群，这就证明了第二个断言。

\begin{enumerate}
\def\labelenumi{(\roman{enumi})}
\setcounter{enumi}{1}
\tightlist
\item
  说 \(L \subset K^{z}\) 即意味着 \(H \supset G^{z}\)，因此当 \([K':K]\) 有限时，(ii) 的断言是第 2 节命题 4 (ii) 与 (iii) 的特殊情形。一般情形下的论证与命题 6 的证明中的相同。
\end{enumerate}

\subsection{第二存在定理}

定理 3. 设 A 为整闭整环，A' 为包含 A 且在 A 上整的环。假设 A 中在 A' 里是零因子的元素只有 0。设 p、q 为 A 的两个素理想，满足 q ⊃ p，并设 q' 为 A' 的位于 q 之上的一个素理想。那么存在 A' 的位于 p 之上的素理想 p'，使 q' ⊃ p'。

先设 A' 是整环。设 K、K' 分别为 A 与 A' 的分式域；设 K'' 为 K' 的代数闭包，A'' 为 A 在 K'' 中的整闭包；于是 \(A \subset A' \subset A''\)。设 \(p''\) 为 A'' 的位于 p 之上的素理想（第 1 节定理 1），\(q''\) 为 A'' 的位于 q 之上且满足 \(p'' \subset q''\) 的素理想（第 1 节定理 1 的推论 2），最后设 \(q_{1}''\) 为 A'' 的位于 \(q'\) 之上的素理想（第 1 节定理 1）。由第 3 节命题 6 (i)，存在 K'' 的一个 K-自同构 \(\sigma\) 使 \(\sigma \cdot q'' = q_{1}''\)。于是 \(\sigma \cdot p''\) 是 A'' 的位于 p 之上且满足 \(\sigma \cdot p'' \subset q_{1}''\) 的素理想，从而 \(p' = A' \cap \sigma \cdot p''\) 是 A' 的位于 p 之上且包含于 \(A' \cap q_{1}'' = q'\) 的素理想。

现在转到一般情形。由于 A 是整环且 \(q'\) 是素理想，子集 \(A - \{0\}\) 与 \(A' - q'\)（皆为 \(A'\) 的子集）都是乘性的；于是它们的积 \(S = (A - \{0\})(A' - q')\) 是 \(A'\) 的一个不含 0 的乘性子集，因为 A 的非零元素都不是 \(A'\) 中的零因子。于是存在（第 II 章 §2 第 5 节命题 11 的推论 2）一个素理想 \(m'\)（属于 \(A'\)），它与 S 不相交，换言之，它满足 \(m' \subset q'\) 且 \(m' \cap A = 0\)。设 h 为典范同态 \(A' \to A'/m'\)。h 在 A 上的限制是单射，因而 \(h(A)\) 是整闭的。由于 \(m' \subset q'\)，\(h(q')\) 是 \(A'/m'\) 的位于 \(h(q)\) 之上的素理想；因为 \(A'/m'\) 是整环，证明的第一部分表明存在一个素理想 \(n'\)（属于 \(A'/m'\)），使 \(n' \cap h(A) = h(p)\) 且 \(h(q') \supset n'\)。理想 \(p' = h^{-1}(n')\) 是 \(A'\) 的素理想，且 \(q' \supset p'\)，因为 \(q'\) 包含 h 的核。由于 \(n' \supset h(p)\)，有 \(p' \supset p\)。最后，对 \(x \in p' \cap A\)，有 \(h(x) \in n' \cap h(A) = h(p)\)，又因 h 在 A 上的限制是单射，故 \(x \in p\)；因此 \(p' \cap A = p\)。

推论. 在定理 3 关于 A 与 A' 的假设下，设 p 为 A 的一个素理想。A' 的位于 p 之上的素理想就是 A' 中包含 pA' 的素理想所成之集 E 的极小元。

A' 的位于 p 之上的素理想依据第 1 节命题 1 的推论 1 是 E 中的极小元。反之，设 q' 为 E 的一个极小元。由于 q' ∩ A ⊃ p，定理 3 表明存在 A' 的位于 p 之上的素理想 p' 使 q' ⊃ p'。由于 p' 包含 pA'，对 q' 所作的假设蕴含 q' = p'，从而 q' 位于 p 之上。

* 设 V、V' 为两个仿射代数簇，f 为 V' 到 V 的一个态射，使得 \(f(V')\) 在 V 中稠密。设 A（相应地 A'）为 V（相应地 V'）上正则函数的环；给定了 f 之后，我们可以把 A 等同于 A' 的一个子环；设 A' 在 A 上整。第 1 节的定理 1 表明，对 V 的每个不可约子簇 W，存在 \(W'\)，即 \(V'\) 的一个不可约子簇，使 \(f(W')\) 是 W 的稠密子集；特别地，V 的每个点都是 \(V'\) 的某个不可约子簇的像，这表明 f 是满射。类似地，f 在每个不可约子簇 \(W'\)（\(V'\) 的）上的限制把 \(W'\) 映到 V 的一个不可约子簇上。第 1 节定理 1 的推论 2 表明：若 W 与 X 是 V 的两个不可约子簇，满足 \(W \supset X\)，且若 \(W'\) 是 \(V'\) 的一个满足 \(f(W') = W\) 的不可约子簇，则存在不可约子簇 \(X'\)（\(V'\) 的、包含于 \(W'\) 的），使 \(f(X') = X\)。

若 A 是整闭的，则称 V 为正规的。定理 3 表明：若 V 正规，若 W 与 X 是 V 的两个不可约子簇，满足 \(W \supset X\)，且若 \(X'\) 是 \(V'\) 的一个满足 \(f(\mathbf{X}') = \mathbf{X}\) 的不可约子簇，则存在子簇 \(W'\)（在 \(V'\) 中），它包含 \(X'\) 且满足 \(f(W') = W\)。最后，定理 3 的推论表明：若 V 正规且 W 是 V 的一个不可约子簇，则不可约子簇 \(W'\)（属于 \(V'\)、满足 \(f(W') = W\) 者）恰是 \(f(W)\) 的不可约分支。

\section{域上的有限生成代数}

\subsection{正规化引理}

在本节及以下各节中，k 表示一个交换域。

定理 1（正规化引理）. 设 A 为一个有限生成 \(k\)-代数，并设 \(\mathfrak{a}_1 \subset \mathfrak{a}_2 \subset \ldots \subset \mathfrak{a}_p\) 为 A 的理想的递增有限序列，满足 \(p \geqslant 1\) 且 \(\mathfrak{a}_p \neq \mathbf{A}\)。那么存在 A 中元素的一个有限序列 \((x_i)_{1 \leqslant i \leqslant n}\)，它们在 \(k\) 上代数无关（第 III 章 § 1 第 1 节），并且满足：

\begin{enumerate}
\def\labelenumi{(\alph{enumi})}
\item
  A 在环 \(\mathbf{B} = k[x_1, \ldots, x_n]\) 上是整的。
\item
  对所有 \(j\)（满足 \(1 \leqslant j \leqslant p\)），存在整数的一个递增序列 \((h(j))_{1 \leqslant j \leqslant p}\)，使得对所有 \(j\)，理想 \(\mathfrak{a}_j \cap \mathbf{B}\)（\(\mathbf{B}\) 中的）由 \(x_1, \ldots, x_{h(j)}\) 生成。
\end{enumerate}

首先注意，只需在 A 是多项式代数 \(k[Y_{1},\ldots,Y_{m}]\) 的情形证明本定理即可。因为在一般情形，A 同构于这样一个代数 \(A'\) 除以理想 \(a_{0}'\) 所得的商；用 \(a_{j}'\) 记 \(a_{j}\) 在 \(A'\) 中的原像，并设 \(x_{i}'(1\leqslant i\leqslant r)\) 是 \(A'\) 的元素，它们对环 \(A'\) 与理想的递增序列 \(a_{0}'\subset a_{1}'\subset\cdots\subset a_{p}'\) 满足命题的条件。于是，像 \(x_{i}\)（即 \(x_{i}'\) 在 A 中的像，其中 \(i>h(0)\)）满足所要的条件；这对条件 (b) 是显然的，而对条件 (a) 则由 §1 第 1 节命题 2 得出；最后，若诸 \(x_{i}(h(0)+1\leqslant i\leqslant r)\) 在 k 上不是代数无关的，就会存在一个非零多项式

\[
\mathrm{Q} \in k [ \mathrm{X} _ {h (0) + 1}, \dots , \mathrm{X} _ {r} ]
\]

使得 \(\mathrm{Q}(x_{h(0)+1}^{\prime},\ldots,x_{r}^{\prime})\in\mathfrak{a}_{0}^{\prime}\cap\mathbf{B}^{\prime}\)，其中记 \(B^{\prime}=k[x_{1}^{\prime},\ldots,x_{r}^{\prime}]\)；但按假设，\(a_{0}^{\prime}\cap B^{\prime}\) 的每个元素都可以唯一地写成系数在 k 中的、诸 \(x_{j}^{\prime}\ (1\leqslant j\leqslant r)\) 的多项式，其每个单项式都至少含有某个 \(x_{j}^{\prime}\)（其中 \(1\leqslant j\leqslant h(0)\)）；于是得到矛盾，这就证明了我们的断言。

因此，在证明的余下部分我们将设 \(A = k[Y_{1}, \ldots, Y_{m}]\)，并对 p 作归纳论证。

\begin{enumerate}
\def\labelenumi{(\Alph{enumi})}
\tightlist
\item
  p = 1。我们区分两种情形：
\end{enumerate}

\textbf{(A1) 理想 \(\mathfrak{a}_1\) 是由一个元素 \(x_1 \notin k\) 生成的主理想。}\par

\(x_{1} = \mathrm{P}(\mathrm{Y}_{1},\ldots ,\mathrm{Y}_{m})\)，其中 \(\mathbf{P}\) 是非常数多项式。我们将看到，适当选取整数 \(r_i > 0\)，环 A 在 \(\mathbf{B} = k[x_1,x_2,\dots ,x_m]\) 上是整的，其中

\[
x _ {i} = \mathrm{Y} _ {i} - \mathrm{Y} _ {1} ^ {r _ {i}} \quad (2 \leqslant i \leqslant m).\tag{1}
\]

为此只需选取诸 \(r_i\) 使 \(Y_1\) 在 \(B\) 上是整的（§ 1 第 1 节命题 4）。现在，有关系式

\[
\mathrm{P} \left(\mathrm{Y} _ {1}, x _ {2} + \mathrm{Y} _ {1} ^ {r _ {2}}, \dots , x _ {m} + \mathrm{Y} _ {1} ^ {r _ {m}}\right) - x _ {1} = 0.\tag{2}
\]

设 \(\mathbf{P} = \sum_{\mathbf{p}} a_{\mathbf{p}} \mathbf{Y}^{\mathbf{p}}\)，其中 \(\mathbf{p} = (p_1, \ldots, p_m)\)，诸 \(p_i\) 是整数 \(\geqslant 0\)，\(\mathbf{Y}^{\mathbf{p}} = \mathbf{Y}_1^{p_1} \ldots \mathbf{Y}_m^{p_m}\)，诸 \(a_{\mathbf{p}}\) 皆 \(\neq 0\) 且属于 \(k\)，并且至少有一个指标 \(\mathbf{p}\) 异于 \((0, \ldots, 0)\)；关系式 (2) 可以写成

\[
\sum_ {\mathbf {p}} a _ {\mathbf {p}} \mathrm{Y} _ {1} ^ {p _ {1}} (x _ {2} + \mathrm{Y} _ {1} ^ {r _ {2}}) ^ {p _ {2}} \dots (x _ {m} + \mathrm{Y} _ {1} ^ {r _ {m}}) ^ {p _ {m}} - x _ {1} = 0.\tag{3}
\]

记 \(f(\mathbf{p}) = p_{1} + r_{2}p_{2} + \cdots + r_{m}p_{m}\)，并设诸 \(r_{i}\) 选取成使诸 \(f(\mathbf{p})\) 互不相同（例如取 \(r_{i} = h^{i}\) 即可，其中 h 是一个严格大于所有 \(p_{j}\) 的整数（《集合论》第 III 章 §5 第 7 节命题 8））。那么将存在唯一的组 \(\mathbf{p} = (p_{1}, \ldots, p_{m})\) 使 \(f(\mathbf{p})\) 达到最大，且关系式 (3) 可以写成

\[
a _ {\mathbf {p}} \mathrm{Y} _ {1} ^ {f (\mathbf {p})} + \sum_ {j <   f (\mathbf {p})} \mathrm{Q} _ {j} (x _ {1}, \dots , x _ {m}) \mathrm{Y} _ {1} ^ {j} = 0\tag{4}
\]

其中诸 \(Q_{j}\) 是 \(k[Y_{1},\ldots,Y_{m}]\) 中的多项式；由于 \(a_{p} \neq 0\) 在 k 中可逆，(4) 当然是系数在 B 中的一个整依赖方程，由此即得我们的断言。

于是分式域 \(k(\mathrm{Y}_1, \ldots, \mathrm{Y}_m)\)（即 \(\mathbf{A}\) 的分式域）在分式域 \(k(x_1, \ldots, x_m)\)（即 \(\mathbf{B}\) 的分式域）上是代数的，这证明（《代数》第 V 章 § 5 第 3 节定理 4）诸 \(x_i\)（\(1 \leqslant i \leqslant m\)）是代数无关的。此外，\(a_1 \cap B = Bx_1\)；因为每个元素 \(z \in a_1 \cap B\) 都可写成 \(z = x_1z'\)，其中 \(z' \in A \cap k(x_1, \ldots, x_m)\)；而由于 B 是整闭的，有 \(A \cap k(x_1, \ldots, x_m) = k[x_1, \ldots, x_m] = B\)（§ 1 第 3 节命题 13 的推论 2）；因此 \(z' \in B\)，这就完成了此情形中性质 (a) 与 (b) 的证明。

\textbf{(A2) 一般情形 \((p = 1)\)。}\par

我们对 \(m\) 作归纳论证，\(m = 0\) 的情形是平凡的。显然可设 \(a_1 \neq 0\)（否则可取 \(x_i = Y_i\)（对 \(1 \leqslant i \leqslant m\)）且取 \(h(1) = 0\)）。设 \(x_1\) 为 \(a_1\) 的一个非零元素；由 (A1)，存在 \(t_2, \ldots, t_m\) 使 \(x_1, t_2, \ldots, t_m\) 在 \(k\) 上代数无关，A 在 \(C = k[x_1, t_2, \ldots, t_m]\) 上是整的，且 \(x_1A \cap C = x_1C\)。由归纳假设，存在元素 \(x_2, \ldots, x_m\)（属于 \(k[t_2, \ldots, t_m]\)）与一个整数 \(h\)，使得 \(k[t_2, \ldots, t_m]\) 在 \(B' = k[x_2, \ldots, x_m]\) 上是整的，\(x_2, \ldots, x_m\) 在 \(k\) 上代数无关，且理想 \(a_1 \cap B'\) 由 \(x_2, \ldots, x_h\) 生成。于是 C 在 \(B = k[x_1, x_2, \ldots, x_m]\) 上是整的（\(\S\) 1 第 1 节命题 5 的推论），从而 A 亦然（\(\S\) 1 第 1 节命题 6）；与情形 (A1) 中同样的论证表明 \(x_1, \ldots, x_m\) 在 \(k\) 上代数无关；最后，由于 \(x_1 \in a_1\) 且 \(B = B'[x_1]\)，有 \(a_1 \cap B = Bx_1 + (a_1 \cap B')\)，又由于 \(a_1 \cap B'\)（在 \(B'\) 中）由 \(x_2, \ldots, x_h\) 生成，故 \(a_1 \cap B\)（在 B 中）由 \(x_1, x_2, \ldots, x_h\) 生成。

\textbf{(B) 从 p - 1 到 p 的过渡。}\par

设 \(t_1, \ldots, t_m\) 为 A 的元素，它们对理想的递增序列 \(a_1 \subset \ldots \subset a_{p-1}\) 满足定理的条件，并记 \(r = h(p - 1)\)。由 (A2)，存在元素 \(x_{r+1}, \ldots, x_m\)（属于 \(k[t_{r+1}, \ldots, t_m]\)）与一个整数 \(s\)，使得

\[
\mathbf {C} = k [ t _ {r + 1}, \dots , t _ {m} ]
\]

在 \(\mathbf{B}' = k[x_{r+1},\ldots,x_m]\) 上是整的，\(x_{r+1},\ldots,x_m\) 在 \(k\) 上代数无关，且理想 \(\mathfrak{a}_p\cap\mathbf{B}'\) 由 \(x_{r+1},\ldots,x_s\) 生成。记 \(x_i=t_i\)（对 \(i\leqslant r\)），则所得到的族 \((x_i)_{1\leqslant i\leqslant m}\) 以 \(h(p)=s\) 解决了问题。因为 A 在 \(\mathbf{C}[t_1,\ldots,t_r]=\mathbf{C}[x_1,\ldots,x_r]\) 上是整的，从而 A 在 \(\mathbf{B}=k[x_1,\ldots,x_m]=\mathbf{B}'[x_1,\ldots,x_r]\) 上也是整的，因为 C 在 \(\mathbf{B}'\) 上是整的（§ 1 第 1 节命题 5 的推论与命题 6）；可以像情形 (A1) 中那样证明诸 \(x_i\) 在 \(k\) 上代数无关。另一方面，对 \(j\leqslant p-1\)，理想

\[
a _ {j} \cap k [ x _ {1}, \dots , x _ {r}, t _ {r + 1}, \dots , t _ {m} ]
\]

按假设是 \(x_1, \ldots, x_r, t_{r+1}, \ldots, t_m\) 中的多项式所成之集，这些多项式的每个单项式都含有元素 \(x_1, \ldots, x_{h(j)}\) 之一；由于 \(x_{r+1}, \ldots, x_m\) 是 \(t_{r+1}, \ldots, t_m\) 的系数在 \(k\) 中的多项式，立即可见，\(x_1, \ldots, x_r, x_{r+1}, \ldots, x_m\)（系数在 \(k\) 中）的多项式属于 \(a_j\) 仅当其所有单项式都含有元素 \(x_1, \ldots, x_{h(j)}\) 之一。最后，由于 \(x_1, \ldots, x_r\) 属于 \(a_{p-1}\) 从而也属于 \(a_p, a_p \cap B'[x_1, \ldots, x_r]\)，它由 \(x_1, \ldots, x_r\) 的、系数在 \(B'\) 中且常数项属于 \(a_p \cap B'\) 的多项式组成；因此这个理想由 \(x_1, \ldots, x_r, x_{r+1}, \ldots, x_t\) 生成。

推论 1. 设 A 为整环，B 为一个包含 A 作为子环的有限生成 A-代数。那么存在 A 的元素 \(s \neq 0\) 与 B 的一个子代数 B'，它同构于多项式代数 \(\mathrm{A}[\mathrm{Y}_{1}, \ldots, \mathrm{Y}_{n}]\)，使得 \(\mathrm{B}[s^{-1}]\)（第 II 章 §2 第 1 节）在 \(\mathrm{B}'[s^{-1}]\) 上是整的。

记 \(S = A - \{0\}\)，并令 \(k = S^{-1}A\)，即 \(A\) 的分式域；显然 \(S^{-1}B\) 是一个有限生成 \(k\)-代数，且由于按假设它包含 \(k\)（第 II 章 § 2 第 4 节定理 1），它不退化为 0。于是由定理 1（应用于 \(p = 1\) 与 \(a_1 = 0\)）存在一个有限序列 \((x_i)_{1 \leqslant i \leqslant n}\)（由 \(S^{-1}B\) 的元素组成），它们在 \(k\) 上代数无关，且使 \(S^{-1}B\) 在 \(k[x_1, \ldots, x_n]\) 上是整的。设 \((z_j)_{1 \leqslant j \leqslant m}\) 为 A-代数 \(B\) 的一个生成系；在 \(S^{-1}B\) 中，每个 \(z_j/1\) 都满足一个整依赖方程

\[
(z _ {j} / 1) ^ {q _ {j}} + \sum_ {h <   q _ {j}} \mathrm{P} _ {h j} (x _ {1}, \dots , x _ {n}) (z _ {j} / 1) ^ {h} = 0\tag{5}
\]

其中诸 \(P_{hj}\) 是 \(x_{i}\) 的系数在 k 中的多项式。存在 A 的元素 \(s \neq 0\)，使我们可以写 \(x_{i} = y_{i}/s\)，其中 \(y_{i} \in B\)（对 \(1 \leqslant i \leqslant n\)），并且诸 \(P_{hj}\) 的所有系数都具有 c/s 的形式，其中 \(c \in A\)；最后，如有必要把 s 换成 S 中元素的一个乘积，我们可假定在 B 中

\[
s z _ {j} ^ {q _ {j}} + \sum_ {h <   q _ {j}} Q _ {h j} z _ {j} ^ {h} = 0\tag{6}
\]

其中诸 \(Q_{hj}\) 是 \(y_{1},\ldots,y_{n}\) 的系数在 A 中的多项式；若记 \(z_{j}^{\prime}=sz_{j}\)，则将 (6) 乘以 \(s^{q_{j}-1}\) 即见 \(z_{j}^{\prime}\) 在 \(B^{\prime}=A[y_{1},\ldots,y_{n}]\) 上是整的。我们证明诸 \(y_{i}\) 在 A 上代数无关；若有一个形如 \(\sum_{p}a_{p}y_{1}^{p_{1}}\ldots y_{n}^{p_{n}}=0\) 的关系，其中对所有 p 有 \(a_{p}\in A\)，则可推出 \(\sum_{p}a_{p}^{\prime}x_{1}^{p_{1}}\ldots x_{n}^{p_{n}}=0\)（在 \(S^{-1}B\) 中），其中 \(a_{p}^{\prime}=a_{p}s^{p_{1}+\cdots+p_{n}}\)（在 k 中）；于是按假设对所有 p 有 \(a_{p}^{\prime}=0\)，从而对所有 p 有 \(a_{p}=0\)。此外，在环 \(B[s^{-1}]\) 中，每个 \(z_{j}^{\prime}/1\) 都在 \(B^{\prime}[s^{-1}]\) 上是整的（§1 第 1 节命题 2），且由于 \(z_{j}/1=(z_{j}^{\prime}/1)(1/s)\)（在 \(B[s^{-1}]\) 中），可见诸 \(z_{j}/1\) 都在 \(B^{\prime}[s^{-1}]\) 上是整的，这就完成了证明（§1 第 1 节命题 4）。

推论 2. 设 K 为域，A 为 K 的子环，L 为 A 的分式域。若 K 是有限生成 A-代数，则 {[}K:L{]} 有限，且存在 A 中的 \(a \neq 0\) 使 \(\mathbf{L} = \mathbf{A}[a^{-1}]\)。

由推论 1，存在元素 \(x_{1}, \ldots, x_{n}\)（属于 \(\mathbf{K}\)）与一个元素 \(a \neq 0\)（属于 \(\mathbf{A}\)），使 \(x_{1}, \ldots, x_{n}\) 在 A 上（从而在 L 上）代数无关，且 K 在子环 \(A[x_{1},\ldots,x_{n},a^{-1}]\) 上是整的。于是由 §2 第 1 节引理 2 可知 \(A[x_{1},\ldots,x_{n},a^{-1}]\) 是一个域。但整环 C 上的多项式环 \(C[Y_{1},\ldots,Y_{n}]\) 中仅有的可逆元就是 C 的可逆元；把这个注记应用于 \(C = A[a^{-1}]\)，便知必有 n = 0，且 \(A[a^{-1}]\) 是一个域，按 L 的定义它等于 L。由于 K 在 L 上是整的且是有限生成 L-代数，次数 {[}K:L{]} 有限（§1 第 1 节命题 4）。

推论 3. 设 A 为整环，B 为有限生成 A-代数，b 为 B 中元素，使得 \(zb^n \neq 0\)（对 A 中所有 \(z \neq 0\) 及每个整数 \(n > 0\)）。设 \(\rho: A \to B\) 为典范同态；存在 A 中的 \(a \neq 0\)，使得对每一个同态 \(f\)（从 A 到代数闭域 L），若满足 \(f(a) \neq 0\)，则存在从 B 到 L 的同态 \(g\)，使 \(g(b) \neq 0\) 且 \(f = g \circ \rho\)。

对 \(b\) 所作的假设蕴含：若 \(h\) 是典范同态 \(x \mapsto x / 1\)（从 \(\mathbf{B}\) 到 \(\mathbf{B}[b^{-1}]\)），则同态 \(h \circ \rho\)（从 \(\mathbf{A}\) 到 \(\mathbf{B}[b^{-1}]\)）是单射。于是由推论 1，存在元素 \(a \neq 0\)（属于 \(\mathbf{A}\)）与子环 \(\mathbf{B}'\)（\(\mathbf{B}[b^{-1}]\) 的），使得 \(\mathbf{B}[b^{-1}, a^{-1}]\) 在 \(\mathbf{B}'[a^{-1}]\) 上是整的，且 \(\mathbf{B}'\) 同构于多项式代数 \(\mathbf{A}[\mathbf{Y}_1, \ldots, \mathbf{Y}_n]\)。设 \(f\) 是 \(\mathbf{A}\) 到代数闭域 \(\mathbf{L}\) 的一个满足 \(f(a) \neq 0\) 的同态；存在从 \(\mathbf{A}[\mathbf{Y}_1, \ldots, \mathbf{Y}_n]\) 到 \(\mathbf{L}\) 的延拓 \(f\) 的同态，从而存在一个同态 \(f'\)（从 \(\mathbf{B}'\) 到 \(\mathbf{L}\)），它延拓 \(f\)。由于 \(f'(a) \neq 0\)（在 \(\mathbf{L}\) 中），存在同态 \(f''\)（从 \(\mathbf{B}'[a^{-1}]\) 到 \(\mathbf{L}\)），使得

\[
f ^ {\prime \prime} (x / a ^ {n}) = f ^ {\prime} (x). (f (a)) ^ {- n}
\]

对所有 \(x \in \mathbf{B}'\) 与所有 \(n > 0\)（第 II 章 §2 第 1 节命题 1）。最后，由于 \(\mathbf{B}[b^{-1}, a^{-1}]\) 在 \(\mathbf{B}'[a^{-1}]\) 上是整的，存在同态 \(f'''\)（从 \(\mathbf{B}[b^{-1}, a^{-1}]\) 到 \(\mathbf{L}\)），它延拓 \(f''\)（§2 第 1 节定理 1 的推论 4）。若 \(j: x \mapsto x / 1\) 是 \(\mathbf{B}\) 到 \(\mathbf{B}[b^{-1}, a^{-1}]\) 的典范同态，则 \(g = f''' \circ j\) 解决了问题，因为 \(j(b)\) 在 \(\mathbf{B}[b^{-1}, a^{-1}]\) 中可逆，从而 \(f'''(j(b)) \neq 0\)（在 \(\mathbf{L}\) 中）。

注意，若在推论 3 中设 B 为整环且 A ⊆ B，则对 b 的假设等价于 \(b \neq 0\)。

\subsection{域上有限生成代数的整闭包}

定理 2. 设 A 为有限生成的整 k-代数，K 为其分式域，A' 为 A 在一个域 K'（K 的有限代数扩张）中的整闭包。那么 A' 是有限生成 A-模，并且是有限生成 k-代数。

由定理 1，存在 A 的一个子代数 C，它同构于多项式代数 \(k[\mathbf{X}_1, \ldots, \mathbf{X}_n]\)，且 A 在 C 上是整的；\(\mathbf{A}'\) 显然是 C 在 \(\mathbf{K}'\) 中的整闭包（§ 1 第 1 节命题 6）；因此我们可以把注意力限制到 \(A = k[X_{1}, \ldots, X_{n}]\) 的情形。设 N 为 \(K'\) 的拟 Galois 扩张（在 \(K'\) 的一个代数闭包中）且由 \(K'\) 生成，它是 K 的有限代数扩张（《代数》第 V 章 § 6 第 3 节命题 9 的推论 1）。只需证明 A 在 N 中的整闭包 B 是有限生成 A-模即可，因为 \(A'\) 是 B 的子 A-模，而 A 是 Noether 环（第 III 章 § 2 第 10 节定理 2 的推论 2）。因此我们可以把注意力限制到 \(K'\) 是 K 的拟 Galois 扩张的情形。这时我们知道（《代数》第 V 章 § 10 第 9 节命题 14），\(K'\) 是 K 的一个（有限）纯不可分扩张 \(K''\) 的（有限）Galois 扩张。若 \(A''\) 是 A 在 \(K''\) 中的整闭包，则 \(A'\) 是 \(A''\) 在 \(K'\) 中的整闭包，于是只需证明 \(A''\) 是有限生成 A-模且 \(A'\) 是有限生成 \(A''\)-模。而一旦证明了 \(A''\) 是有限生成 A-模，它就是一个 Noether 整环，且按定义是整闭的；\(A'\) 是有限生成 \(A''\)-模这一事实将由 § 1 第 6 节命题 18 的推论 1 得出。

于是可见，我们可以把注意力限制到如下情形：\(\mathbf{A} = k[\mathbf{X}_1, \ldots, \mathbf{X}_n]\)，且 \(\mathbf{K}'\) 是 \(\mathbf{K} = k(\mathbf{X}_1, \ldots, \mathbf{X}_n)\) 的有限纯不可分扩张。这时 \(\mathbf{K}'\) 由元素的一个有限族 \((y_i)_{1 \leqslant i \leqslant m}\) 生成，且存在 \(q\)（\(k\) 的特征指数的一个幂），使 \(y_i^q \in k(\mathbf{X}_1, \ldots, \mathbf{X}_n)\)。设 \(c_j\)（\(1 \leqslant j \leqslant r\)）为 \(\mathbf{X}_1, \ldots, \mathbf{X}_n\) 的、等于 \(y_i^q\)（\(1 \leqslant i \leqslant m\)）的有理函数之分子与分母的系数。那么 \(\mathbf{K}'\) 包含于扩张 \(\mathbf{L} = k'(\mathbf{X}_1^{q^{-1}}, \ldots, \mathbf{X}_n^{q^{-1}})\)，其中 \(k' = k(c_1^{q^{-1}}, \ldots, c_n^{q^{-1}})\)（我们在 \(\mathbf{K}'\) 的一个代数闭包中），且 \(\mathbf{A}'\) 包含于代数闭包 \(\mathbf{B}'\)（即 \(\mathbf{A}\) 在 \(\mathbf{L}'\) 中的代数闭包）。现在，\(k'\) 在 \(k\) 上是代数的，因此 \(\mathbf{C}' = k'[\mathbf{X}_1, \ldots, \mathbf{X}_n]\) 在 \(\mathbf{A}\) 上是整的（\(\S 1\) 第 1 节命题 5）；由于 \(k'[\mathbf{X}_1^{q^{-1}}, \ldots, \mathbf{X}_n^{q^{-1}}]\) 是整闭的（\(\S 1\) 第 3 节命题 13 的推论 2），可见这个环是 \(\mathbf{C}'\) 在 \(\mathbf{L}'\) 中的整闭包，从而也是 \(\mathbf{A}\) 的整闭包（\(\S 1\) 第 1 节命题 6），换言之，\(\mathbf{B}' = k'[\mathbf{X}_1^{q^{-1}}, \ldots, \mathbf{X}_n^{q^{-1}}]\)。现在显然 \(\mathbf{B}'\) 是有限生成 \(\mathbf{C}'\)-模（\(\S 1\) 第 1 节命题 4），且由于 \(k'\) 是 \(k\) 的有限扩张，\(\mathbf{C}'\) 是有限生成 A-模，从而 \(\mathbf{B}'\) 是有限生成 A-模；由于 \(\mathbf{A}\) 是 Noether 的且 \(\mathbf{A}' \subset \mathbf{B}'\)，故 \(\mathbf{A}'\) 是有限生成 A-模。

\subsection{零点定理}

命题 1. 设 A 为域 \(k\) 上的一个有限生成代数，\(L\) 为 \(k\) 的代数闭包。

\begin{enumerate}
\def\labelenumi{(\roman{enumi})}
\item
  若 \(\mathbf{A} \neq \{0\}\)，则存在一个 \(k\)-同态（从 \(\mathbf{A}\) 到 \(\mathbf{L}\)）。
\item
  设 \(f_{1}, f_{2}\) 为两个 \(k\)-同态（从 \(\mathbf{A}\) 到 \(\mathbf{L}\)）。要使 \(f_{1}\) 与 \(f_{2}\) 有相同的核，必须且只需存在一个 \(k\)-自同构 \(s\)（\(\mathbf{L}\) 的），使 \(f_{2} = s \circ f_{1}\)。
\item
  设 \(\mathfrak{a}\) 为 \(\mathbf{A}\) 的一个理想。要使 \(\mathfrak{a}\) 是极大的，必须且只需它是一个 \(k\)-同态（从 \(\mathbf{A}\) 到 \(\mathbf{L}\)）的核。
\item
  要使元素 \(x\)（\(\mathbf{A}\) 的）满足 \(f(x) = 0\)（对每个 \(k\)-同态 \(f\)，从 \(\mathbf{A}\) 到 \(\mathbf{L}\)），必须且只需 \(x\) 是幂零的。
\end{enumerate}

断言 (i) 由第 1 节定理 1 的推论 3 得出，应用时把 A 换成 k，B 换成 A，b 换成 B 的单位元，f 换成 k 到 L 的典范单射。

若 \(f\) 是一个 \(k\)-同态（从 \(\mathbf{A}\) 到 \(\mathbf{L}\)），则 \(f(\mathbf{A})\) 是 \(\mathbf{L}\) 的一个包含 \(k\) 的子环；由于 \(\mathbf{L}\) 是 \(k\) 的代数扩张，\(f(\mathbf{A})\) 是一个域（《代数》第 V 章 § 3 第 2 节命题 3），而若 \(\mathfrak{a}\) 是 \(f\) 的核，则 \(\mathbf{A} / \mathfrak{a}\)（同构于 \(f(\mathbf{A})\)）因此是一个域，这就证明 \(\mathfrak{a}\) 是极大的。反之，若 \(\mathfrak{a}\) 是 \(\mathbf{A}\) 的一个极大理想，则由 (i) 可知存在一个 \(k\)-同态（从 \(\mathbf{A} / \mathfrak{a}\) 到 \(\mathbf{L}\)），从而存在一个 \(k\)-同态（\(\mathbf{A}\) 到 \(\mathbf{L}\)），其核 \(\mathfrak{b}\) 包含 \(\mathfrak{a}\)；但由于 \(\mathfrak{a}\) 极大，\(\mathfrak{b} = \mathfrak{a}\)；这就证明了 (iii)。

现在证明 (ii)。若 \(s\) 是一个 \(k\)-自同构（\(L\) 的），满足 \(f_2 = s \circ f_1\)，则显然 \(f_1\) 与 \(f_2\) 有相同的核。反之，设 \(f_1\) 与 \(f_2\) 有相同的核；那么存在一个 \(k\)-同构 \(s_0\)（从域 \(f_1(A)\) 到域 \(f_2(A)\) 上），使 \(f_2 = s_0 \circ f_1\)；但由《代数》第 V 章 § 6 第 3 节命题 7，\(s_0\) 可延拓为一个 \(k\)-自同构 \(s\)（\(L\) 的），从而 \(f_2 = s \circ f_1\)。

最后，若 \(x \in \mathbf{A}\) 满足 \(x^n = 0\)，则对每个 \(k\)-同态 \(f\)（从 \(\mathbf{A}\) 到 \(\mathbf{L}\)），有 \((f(x))^n = f(x^n) = 0\)，从而 \(f(x) = 0\)，因为 \(\mathbf{L}\) 是域。反之，设 \(x \in \mathbf{A}\) 不是幂零的；那么 \(\mathbf{A}[x^{-1}]\) 是一个不退化为 0 的有限生成 A-代数（从而是有限生成 \(k\)-代数）（第 II 章 §2 第 1 节注 3），于是由 (i) 存在一个 \(k\)-同态 \(g\)（从 \(\mathbf{A}[x^{-1}]\) 到 \(\mathbf{L}\)）。若 \(j: \mathbf{A} \to \mathbf{A}[x^{-1}]\) 是典范同态，则 \(f = g \circ j\) 是一个 \(k\)-同态（从 \(\mathbf{A}\) 到 \(\mathbf{L}\)），且 \(f(x)g(1/x) = g(x/1)g(1/x) = g(1) = 1\)，故 \(f(x) \neq 0\)。

设 \(k\) 为域，\(\mathbf{L}\) 为 \(k\) 的一个扩域；元素 \(\mathbf{x} = (x_1, \ldots, x_n)\)（\(\mathbf{L}^n\) 中的）称为 \(\mathbf{L}^n\) 中的一个零点，即理想 \(\mathfrak{r}\)（多项式环 \(k[\mathbf{X}_1, \ldots, \mathbf{X}_n]\) 的理想）的零点，如果

\[
\mathrm{P} (\mathbf {x}) = \mathrm{P} (x _ {1}, \dots , x _ {n}) = 0
\]

对所有 \(\mathbf{P} \in \mathfrak{r}\)。

引理 1. 设 A 为域 \(k\) 上的一个有限生成代数，\((a_i)_{1 \leqslant i \leqslant n}\) 为它的一个生成系，\(\mathfrak{r}\) 为诸 \(a_i\) 之间的、系数在 \(k\) 中的代数关系理想（《代数》第 IV 章 §2 第 1 节）。对每个扩域 \(\mathbf{L}\)（\(k\) 的），映射 \(f \mapsto (f(a_i))_{1 \leqslant i \leqslant n}\) 是从 A 到 L 的 \(k\)-同态所成之集到 \(\mathfrak{r}\) 在 \(L^n\) 中的零点所成之集上的一个双射。

存在唯一的一个 \(k\)-代数同态 \(h\)，从 \(k[\mathbf{X}_1, \ldots, \mathbf{X}_n]\) 到 \(\mathbf{A}\) 上，使 \(h(\mathbf{X}_i) = a_i\)（对 \(1 \leqslant i \leqslant n\)），且按定义 \(\mathfrak{r}\) 是 \(h\) 的核。映射 \(f \mapsto f \circ h\) 是诸 \(k\)-同态（从 \(A\) 到 \(L\)）所成之集到诸 \(k\)-同态（从 \(k[X_1, \ldots, X_n]\) 到 \(L\)、在 \(r\) 上为零）所成之集上的一个双射。对每个多项式 \(P \in k[X_1, \ldots, X_n]\) 与每个元素 \(\mathbf{x} = (x_1, \ldots, x_n) \in L^n\)，记 \(h_{\mathbf{x}}(P) = P(\mathbf{x})\)；于是映射 \(\mathbf{x} \mapsto h_{\mathbf{x}}\) 是从 \(L^n\) 到诸 \(k\)-同态（从 \(k[X_1, \ldots, X_n]\) 到 \(L\)）所成之集上的一个双射（这样的同态由它在诸 \(X_i\)（\(1 \leqslant i \leqslant n\)）处的值确定）；说 \(h_{\mathbf{x}}\) 在 \(r\) 上为零，就是说 \(\mathbf{x}\) 是 \(r\) 在 \(L^n\) 中的一个零点，由此即得引理。

若把命题 1 应用于代数 \(\mathbf{A} = k[\mathbf{X}_1, \ldots, \mathbf{X}_n] / \mathfrak{r}\)，其中 \(\mathfrak{r}\) 是 \(k[\mathbf{X}_1, \ldots, \mathbf{X}_n]\) 的一个异于整个环的理想，则由引理 1 得到下述命题：

命题 2（Hilbert 零点定理）. 设 \(k\) 为域，\(L\) 为 \(k\) 的一个代数闭包。

\begin{enumerate}
\def\labelenumi{(\roman{enumi})}
\item
  每个理想 \(\mathfrak{r}\)（\(k[\mathbf{X}_1, \ldots, \mathbf{X}_n]\) 的），若不含 1，则在 \(\mathbf{L}^n\) 中至少有一个零点。
\item
  设 \(\mathbf{x} = (x_1, \ldots, x_n)\)、\(\mathbf{y} = (y_1, \ldots, y_n)\) 为 \(\mathbf{L}^n\) 的两个元素；要使 \(k[\mathbf{X}_1, \ldots, \mathbf{X}_n]\) 中在 \(\mathbf{x}\) 处为零的多项式所成之集与 \(k[\mathbf{X}_1, \ldots, \mathbf{X}_n]\) 中在 \(\mathbf{y}\) 处为零的多项式所成之集相同，必须且只需存在一个 \(k\)-自同构 \(s\)（\(\mathbf{L}\) 的），使 \(y_i = s(x_i)\)（对 \(1 \leqslant i \leqslant n\)）。
\item
  要使理想 \(\mathfrak{a}\)（\(k[\mathbf{X}_1, \ldots, \mathbf{X}_n]\) 的）是极大的，必须且只需存在一个 \(\mathbf{x}\)（在 \(\mathbf{L}^n\) 中），使 \(\mathfrak{a}\) 是 \(k[\mathbf{X}_1, \ldots, \mathbf{X}_n]\) 中在 \(\mathbf{x}\) 处为零的多项式所成之集。
\item
  要使多项式 \(\mathbf{Q}\)（\(k[\mathbf{X}_1, \ldots, \mathbf{X}_n]\) 的）在 \(\mathbf{L}^n of\) 中、理想 \(\mathfrak{r}\)（\(k[\mathbf{X}_1, \ldots, \mathbf{X}_n]\) 的）的零点所成之集上为零，必须且只需存在一个整数 \(m > 0\) 使 \(\mathbf{Q}^m \in \mathfrak{r}\)。
\end{enumerate}

\subsection{Jacobson 环}

定义 1. 若环 \(\mathbf{A}\) 的每个素理想都是一族极大理想的交，则称 \(\mathbf{A}\) 为 Jacobson 环。

\textbf{例子}\par

\begin{enumerate}
\def\labelenumi{(\arabic{enumi})}
\item
  每个域都是 Jacobson 环。
\item
  环 Z 是 Jacobson 环，因为唯一不是极大理想的素理想 (0) 是 Z 的诸极大理想 \((p)\) 的交，其中 p 取遍素数集（参看命题 4）。
\item
  设 A 为 Jacobson 环，a 为 A 的理想。那么 A/a 是 Jacobson 环，因为 A/a 的理想都具有 b/a 的形式，其中 b 是 A 的包含 a 的理想，且 b/a 是素理想（相应地极大理想）当且仅当 b 是。
\end{enumerate}

命题 3. 要使环 A 是 Jacobson 环，必须且只需对 A 的每个理想 a，A/a 的 Jacobson 根等于其幂零根（第 II 章 § 2 第 6 节）。

A/a 的 Jacobson 根（相应地幂零根）是 A/a 的极大理想（相应地素理想）的交（《代数》第 VIII 章 § 6 第 3 节定义 3 及《交换代数》第 II 章 § 2 第 6 节命题 13）。因此所述条件意味着：对 A 的每个理想 a，包含 a 的素理想之交等于包含 a 的极大理想之交。若 A 是 Jacobson 环，此条件显然对 A 的每个理想 a 成立；反之，若它对 A 的每个素理想成立，则按定义 A 是 Jacobson 环。

推论. 设 A 为 Jacobson 环；对 A 的每个理想 a，a 的根是 A 的包含 a 的极大理想的交。

只需注意 A/a 是 Jacobson 环即可。

命题 4. 设 A 为一个主理想整环，\((p_{\lambda})_{\lambda \in L}\) 为 A 的极端元的一个代表系（《代数》第 VII 章 § 1 第 3 节定义 2）。要使 A 是 Jacobson 环，必须且只需 L 是无限的。

A 的极大理想是诸 \(Ap_{\lambda}\)（同处第 2 节命题 2）。若 L 有限，它们的交是理想 Ax，其中 \(x = \prod_{\lambda \in L} p_{\lambda}\)（同处），因而异于 (0)；另一方面，若 L 无限，则诸 \(Ap_{\lambda}\) 的交是 (0)，因为 A 的每个 \(\neq 0\) 的元素只被有限多个极端元整除（同处第 3 节定理 2）。于是命题由 (0) 是 A 中唯一不是极大理想的素理想这一事实得出（《代数》第 VI 章 §1 第 13 节命题 14 (DIV)）。

命题 5. 设 A 为环，B 为一个在 A 上整的 A-代数。若 A 是 Jacobson 环，则 B 亦然。

把 A 换成它在 B 中的典范像，我们可设 \(A \subset B\)。设 \(p'\) 为 B 的一个素理想，并令 \(p = A \cap p'\)。按假设存在 A 的一族极大理想 \((m_{\lambda})_{\lambda \in L}\)，其交等于 p。对每个 \(\lambda \in L\)，存在 B 的一个极大理想 \(m_{\lambda}'\)，它位于 \(m_{\lambda}\) 之上且包含 \(p'\)（§ 2 第 1 节命题 1 及定理 1 的推论 2）。若记 \(q' = \bigcap_{\lambda \in L} m_{\lambda}'\)，则 \(q' \cap A = \bigcap_{\lambda \in L} m_{\lambda} = p\) 且 \(q' \supset p'\)，从而 \(q' = p'\)（§ 2 第 1 节命题 1 的推论 1）。

定理 3. 设 A 为 Jacobson 环，B 为有限生成 A-代数，\(\rho: \mathbf{A} \to \mathbf{B}\) 为典范同态。那么：

\begin{enumerate}
\def\labelenumi{(\roman{enumi})}
\item
  B 是 Jacobson 环。
\item
  对每个极大理想 \(\mathfrak{m}'\)（\(\mathbf{B}\) 的），\(\mathfrak{m} = \rho^{-1} (\mathfrak{m}')\) 是 \(\mathbf{A}\) 的极大理想，且 \(\mathbf{B} / \mathfrak{m}'\) 是 \(\mathbf{A} / \mathfrak{m}\) 的有限代数扩张。
\end{enumerate}

设 \(\mathfrak{p}'\) 为 B 的一个素理想，\(\mathfrak{p} = \rho^{-1} (\mathfrak{p}')\)。设 \(v\) 为一个 \(\neq 0\) 的元素（在 \(\mathrm{B / p'}\) 中）。

由于 \(B/p'\) 是一个有限生成的整 \((A/p)\)-代数，且典范同态 \(\phi: A/p \to B/p'\) 是单射，故存在一个元素 \(u \neq 0\)（属于 \(A/p\)），使得：对每个同态 \(f\)（从 \(A/p\) 到代数闭域 \(L\)，其核不含 \(u\)），存在一个同态 \(g\)（从 \(B/p'\) 到 \(L\)），其核不含 \(v\) 且满足 \(f = g \circ \phi\)（第 1 节定理 1 的推论 3）。由于 \(A\) 是 Jacobson 环，存在一个极大理想 \(m\)（属于 \(A\)），它包含 \(p\) 且使 \(u \notin m/p\)。取 \(L\) 为 \(A/m\) 的一个代数闭包，取 \(f\) 为典范同态 \(A/p \to L\)；设

\[
g \colon \mathrm{B} / \mathfrak {p} ^ {\prime} \rightarrow \mathrm{L}
\]

为一个满足 \(f = g \circ \phi\) 且 \(g(v) \neq 0\) 的同态。那么

\[
\mathbf {A} / \mathfrak {m} \subset g (\mathbf {B} / \mathfrak {p} ^ {\prime}) \subset \mathbf {L},
\]

从而 \(g(\mathbf{B} / \mathfrak{p}')\) 是 \(\mathbf{K}\) 的一个子域（《代数》第 V 章 § 3 第 2 节命题 3），因此 \(g\) 的核是 \(\mathbf{B} / \mathfrak{p}'\) 的一个不含 \(v\) 的极大理想。于是可见 \(\mathbf{B} / \mathfrak{p}'\) 的诸极大理想之交化为 0，这就证明 \(\mathbf{B}\) 是 Jacobson 环。此外，若 \(\mathfrak{p}'\) 极大，则 \(g\) 必为单射，从而 \(\mathfrak{p} = \mathfrak{m}\) 极大；最后，此时 \(\mathbf{B} / \mathfrak{p}'\) 是域 \(\mathbf{A} / \mathfrak{m}\) 上的有限生成代数，从而是 \(\mathbf{A} / \mathfrak{m}\) 的有限扩张（第 1 节定理 1 的推论 2）。

推论 1. Z 上的每个有限生成代数 A 都是 Jacobson 环；要使 A 的素理想 p 是极大的，必须且只需环 A/p 有限。

若整环 A/p 有限，则它是一个域，因为对 A/p 中每个 \(u \neq 0\)，A/p 到自身的映射 \(v \mapsto uv\) 是单射，从而由于 A/p 有限而是双射。反之，对 A 的每个极大理想 m，m 在 Z 中的原像是一个极大理想 (p)，且由定理 3，A/m 在素域 \(\mathbf{Z}/(\mathfrak{p}) = \mathbf{F}_{\mathfrak{p}}\) 上有限。

推论 2. 设 \((\mathbf{P}_{\lambda})_{\lambda \in \mathbf{L}}\) 为 \(\mathbf{Z}[\mathbf{X}_1,\ldots ,\mathbf{X}_n]\) 中的一个多项式族，并设 \(\mathbf{Q}\) 为 \(\mathbf{Z}[\mathbf{X}_1,\dots ,\mathbf{X}_n]\) 中的一个多项式，使得：对每个属于一个有限域的元素组 \((x_i)_{1\leqslant i\leqslant n}\)，若 \(\mathrm{P}_{\lambda}(x_1,\ldots ,x_n) = 0\)（对所有 \(\lambda\)），则也有 \(\mathrm{Q}(x_1,\ldots ,x_n) = 0\)。那么，若 \(\mathfrak{a}\) 是 \(\mathbf{Z}[\mathbf{X}_1,\ldots ,\mathbf{X}_n]\) 中由诸 \(\mathrm{P}_{\lambda}\) 生成的理想，则存在整数 \(m > 0\) 使 \(\mathbf{Q}^m\in \mathfrak{a}\)。此外，对每个约化环 R 与 R 中元素的每个组 \((y_i)_{1\leqslant i\leqslant n}\)，若 \(\mathrm{P}_{\lambda}(y_1,\ldots ,y_n) = 0\)（对所有 \(\lambda\)），则也有 \(\mathrm{Q}(y_1,\ldots ,y_n) = 0\)。

第二个断言由第一个断言得出，因为 \(\mathbf{Z}[\mathrm{X}_1,\ldots ,\mathrm{X}_n]\) 的由多项式 \(\mathbf{P}\)（满足 \(\mathrm{P}(y_1,\dots,y_n) = 0\)）组成的理想包含 a。为证明第一个断言，只需注意：对 \(\mathbf{A} = \mathbf{Z}[\mathrm{X}_1,\dots,\mathrm{X}_n]\) 的每个包含 a 的极大理想 m，A/m 是一个有限域（推论 1），且假设蕴含 Q 在 A/m 中的典范像为零；于是 Q 属于 A 的包含 a 的诸极大理想之交，即 a 的根（命题 3 的推论）。

推论 3. 设 A 为 Jacobson 环。若存在一个包含 A 且是域的有限生成 A-代数 B，则 A 是域而 B 是 A 的代数扩张。

只需对 \(\mathfrak{m}' = (0)\) 应用定理 3 (ii)。


\exercisesection{1}

\begin{enumerate}
\def\labelenumi{\arabic{enumi}.}
\item
  设 \(d\) 为一个在 \(\mathbf{Z}\) 中不被平方数整除的有理整数。证明：域 \(\mathbf{K} = \mathbf{Q}(\sqrt{d})\) 中在 \(\mathbf{Z}\) 上整的元素形如 \(a + b\sqrt{d}\)（其中 \(a \in \mathbf{Z}\)、\(b \in \mathbf{Z}\)），若 \(d - 1 \not\equiv 0 \pmod{4}\)；而形如 \((a + b\sqrt{d})/2\)（其中 \(a \in \mathbf{Z}\)、\(b \in \mathbf{Z}\)，且 \(a\) 与 \(b\) 同为偶数或 \(a\) 与 \(b\) 同为奇数），若 \(d - 1 \equiv 0 \pmod{4}\)（参见《代数》第七章 § 1 习题 8）。由此推出 \(\mathbf{A} = \mathbf{Z}[\sqrt{5}]\) 不是整闭的，并给出 \(\mathbf{Q}[\sqrt{5}]\) 中一个在 \(\mathbf{A}\) 上整的元素的例子，但它在 \(\mathbf{A}\) 的分式域上的极小多项式的系数不属于 \(\mathbf{A}\)。
\item
  设 K 为交换域，A 为多项式代数 \(K[X, Y]\) 中由单项式 \(Y^{k}X^{k+1} (k \geqslant 0)\) 生成的子 K-代数。证明：XY 使得 A{[}XY{]} 包含于一个有限生成的 A-模，但 XY 不在 A 上整。
\item
  给出交换域 K 的扩张的一个无穷序列 \((\mathbf{K}_{n})\) 的例子，其中每个扩张在 K 上有限次，使得 K-代数 \(\prod_{n} K_{n}\) 不在 K 上代数。
\item
  在矩阵环 \(\mathbf{M}_{2}(\mathbf{Q})\) 中给出两个在 Z 上整的元素的例子，但它们的和与积都不在 Z 上整（考虑形如 \(I + N\) 的矩阵，其中 N 幂零）。
\item
  设 A 为交换环，B 为包含 A 的交换环，x 为 B 的一个可逆元。证明 \(A[x] \cap A[x^{-1}]\) 的每个元素都在 A 上整。（若 \(y = a_{0} + \cdots + a_{p}x^{-p} = b_{0} + \cdots + b_{q}x^{q}\)，其中 \(a_{i}\) 与 \(b_{j}\) 属于 A，证明由 1，\(x, \ldots, x^{p+q+1}\) 生成的 B 的子 A-模是一个忠实 A{[}y{]}-模。）
\item
  设 A 为交换环，B 为交换 A-代数；假设 B 的极小素理想个数有限；设 \(p_{i}\)（\(1 \leqslant i \leqslant n\)）为这些理想（注意若 B 是 Noether 环则此假设成立）。元素 \(x \in B\) 在 A 上整，当且仅当它的每个典范像 \(x_{i}\)（在 \(B/p_{i}\) 中）都在 A 上整。（若此条件成立，先证明 x 在 B 的由幂零根 \(N = \bigcap_{i} p_{i}\) 生成的子代数上整，并利用 N 的元素皆幂零这一事实。）
\item
  给出一个约化环 A 的例子，它在其全分式环中整闭，但不是整环之积（取域 K 在无穷积 \(\prod_{n} K_{n}\)（即 K 的适当代数扩张之积）中的整闭包）。
\item
  设 A 为交换环，B 为有限生成 A-代数。证明存在一个作为自由 A-模的有限 A-代数 C 以及一个满射 A-同态 \(C \rightarrow B\)。
\item
  证明商整环 \(\mathbf{Z}[\mathbf{X}] / (\mathbf{X}^2 +4)\) 不是整闭的，尽管 \(\mathbf{Z}[\mathbf{X}]\) 是整闭的。
\item
  设 A 为完全整闭整环。证明：对每个集合 I，多项式整环 \(A[X_{i}]_{i \in I}\) 与形式幂级数整环 \(A[[X_{i}]]_{i \in I}\)（《代数》第四章 §5 习题 1）都是完全整闭的。（利用第 4 节命题 4 及下面的引理：对 I 的每个有限子集 J 及所有 \(P \in A[[X_{i}]]_{i \in J}\)，设 \(P_{J}\) 为 \(A[[X_{i}]]_{i \in J}\) 中由 P 出发、把 P 中所有 \(X_{i}\)（指 \(i \notin J\) 者）换为 0 而得的形式幂级数；若 P、Q 是两个形式幂级数，使得 \(P_{J}\) 被 \(Q_{J}\) 整除（在 \(A[[X_{i}]]_{i \in J}\) 中，对所有 J），则 P 在 \(A[[X_{i}]]_{i \in I}\) 中被 Q 整除。为此注意：若 J、\(J'\) 是 I 的两个有限子集，满足 \(J \subset J'\)，且 \(P_{J'} = L'Q_{J'}\) 与 \(P_{J} = LQ_{J}\)，则 L 由 \(L'\) 出发、把 \(L'\) 中所有 \(X_{i}\)（指 \(i \in J' - J\) 者）换为 0 而得到。）
\item
  \begin{enumerate}
  \def\labelenumii{(\alph{enumii})}
  \tightlist
  \item
    设 R 为整环，\((\mathbf{A}_{\alpha})\) 为 R 的一个右定向的子环族，A 为诸 \(A_{\alpha}\) 的并。证明：若每个 \(A_{\alpha}\) 都整闭，则 A 也整闭。
  \end{enumerate}
\end{enumerate}

\begin{enumerate}
\def\labelenumi{(\alph{enumi})}
\setcounter{enumi}{1}
\tightlist
\item
  设 K 为交换域，\(\mathbf{R} = \mathbf{K}(\mathbf{X}, \mathbf{Y})\) 为 K 上两个未定元的有理函数域。对每个整数 n \textgreater{} 0，设 \(A_{n}\) 表示 R 的子环 \(\mathbf{K}[\mathbf{X}, \mathbf{Y}, \mathbf{Y}\mathbf{X}^{-n}]\)。证明诸 \(A_{n}\) 都是完全整闭的，但它们的并 A 不是。（为证形如 \(\mathbf{P}(\mathbf{X}, \mathbf{Y}) \cdot \mathbf{X}^{-ns}\) 的元素（其中 \(P \in K[X, Y]\)）属于 \(A_{n}\)，证明其必要且充分条件是 P 属于 \(S^{s}\)，其中 S 表示 \(K[X, Y]\) 中由 \(X^{n}\) 与 Y 生成的理想。）
\end{enumerate}

¶ * 12. 设 A 为复变量的（取复值的）整函数环。它是一个整环，其分式域是复变量的亚纯函数域。

\begin{enumerate}
\def\labelenumi{(\alph{enumi})}
\item
  证明 \(\mathbf{A}\) 是完全整闭的。
\item
  对所有 \(x \in \mathbf{A}\)，设 \(Z(x)\) 为 \(x\) 的零点集（它是 \(\mathbf{C}\) 的闭子集，当 \(x \neq 0\) 时是离散的（从而可数））。若 \(\mathfrak{a}\) 是 \(\mathbf{A}\) 的一个异于 \(\mathbf{A}\) 与 (0) 的理想，证明诸 \(Z(x)\)（对 \(x \in \mathfrak{a}\)）构成一个滤子基 \(\mathfrak{F}_{\mathfrak{a}}\)。（若 \(x\) 与 \(y\) 是两个没有公共零点的整函数，试利用 Mittag-Leffler 定理证明可以写出 \(1 / xy = u / x + v / y\)，其中 \(u\) 与 \(v\) 属于 \(\mathbf{A}\)，并由此推出 \(Ax + Ay = A\)。）反之，对每个滤子 \(\mathfrak{F}\)（在 \(\mathbf{C}\) 上、且 \(\mathbf{C}\) 的某个闭离散子集属于它），集合 \(\mathfrak{a}(\mathfrak{F})\)（由 \(x \in \mathbf{A}\) 中使 \(Z(x) \in \mathfrak{F}\) 者组成）是 \(\mathbf{A}\) 的一个理想。证明：\(\mathfrak{F}_{\mathfrak{a}(\mathfrak{F})} = \mathfrak{F}\) 对每个滤子 \(\mathfrak{F}\)（\(\mathbf{C}\) 的某个闭离散子集属于它）成立，且有 \(\mathfrak{a}(\mathfrak{F}) \supset \mathfrak{a}\) 对每个理想 \(\mathfrak{a}\)（\(\mathbf{A}\) 的、异于 \(\mathbf{A}\) 与 (0) 的）成立。
\item
  若 \(\mathfrak{p} \neq (0)\) 是 \(\mathbf{A}\) 的素理想，证明 \(\mathfrak{F}_{\mathfrak{p}}\) 是超滤子。（限于 \(\mathfrak{F}_{\mathfrak{p}}\) 中所有集合都是无穷集的情形，证明：若 \(\mathbf{M} \in \mathfrak{F}_{\mathfrak{p}}\) 是两个无公共点的无穷集 \(\mathbf{M}'\)、\(\mathbf{M}''\) 之并，则其中之一属于 \(\mathfrak{F}_{\mathfrak{p}}\)。）反之，对每个超滤子 \(\mathfrak{U}\)（在 \(\mathbf{C}\) 上、且 \(\mathbf{C}\) 的某个闭离散子集属于它），\(\mathfrak{a}(\mathfrak{U})\) 是 \(\mathbf{A}\) 的一个极大理想；试得出其逆。
\item
  设 \(\mathfrak{U}\) 为 \(\mathbf{C}\) 上的一个超滤子，它是在典范单射 \(\mathbf{Z} \to \mathbf{C}\) 下、一个 \(\mathbf{Z}\) 上的非平凡超滤子的像。证明存在非极大的素理想 \(\mathfrak{p}\)（\(\mathbf{A}\) 的）使 \(\mathfrak{F}_{\mathfrak{p}} = \mathfrak{U}\)。（对所有 \(x \in \mathfrak{a}(\mathfrak{U})\) 及所有 \(n \in \mathbf{Z}\)，设 \(\omega_x(n)\) 为 \(x\) 在点 \(n\) 处的阶（\(\omega_x(n) = 0\)，若 \(x(n) \neq 0\)）；考虑集合 \(\mathfrak{p}\)：它由诸 \(x \in \mathfrak{a}(\mathfrak{U})\)（使 \(\lim_{\mathfrak{u}} \omega_x = +\infty\) 者）组成。）
\item
  用 (d) 的记号，设 m 为极大理想 \(\mathfrak{a}(\mathfrak{U})\)。证明整环 \(A_{m}\) 不是完全整闭的（考虑亚纯函数 \(1/(\sin\pi z)\)）。
\end{enumerate}

¶ 13. 设 A 为局部整环，K 为其分式域。考虑下面两个性质：

\begin{enumerate}
\def\labelenumi{(\roman{enumi})}
\item
  A 是 Hensel 的（第三章 § 4 习题 3）。
\item
  对每个有限扩张 \(\mathbf{L}\)（\(\mathbf{K}\) 的），整闭包 \(\mathbf{A}'\)（即 \(\mathbf{A}\) 在 \(\mathbf{L}\) 中的整闭包）是局部环。
\end{enumerate}

证明 (i) 蕴含 (ii)，且当 A 整闭时其逆为真。（为证 (i) 蕴含 (ii)，注意 A' 是一个由有限 A-代数组成的右定向族的并，并利用 Hensel 环的定义。为证当 A 整闭时 (ii) 蕴含 (i)，证明：若 P ∈ A{[}X{]} 是首一不可约多项式，且 f: A → k 是 A 到其剩余域 k 上的典范同态，则 \(\bar{f}(P)\) 在 K{[}X{]} 中不可约，并考虑域 L = K{[}X{]}/(f)。）

若条件 (ii) 满足，则 A 在 K 的任何扩张中的整闭包 A' 都是局部环。考察完备 Hausdorff 局部环的情形；若 A 是完备 Hausdorff Noether 局部环，则每个含于 A' 且作为有限生成 A-模的 A-代数都是完备 Hausdorff 局部环。

\begin{enumerate}
\def\labelenumi{\arabic{enumi}.}
\setcounter{enumi}{13}
\tightlist
\item
  设 A 为完全整闭整环，K 为其分式域。证明对 K 的每个扩张 L，A 在
\end{enumerate}

L 中的整闭包 A' 是完全整闭整环。（化归到 L 是 K 的次数为 m 的有限拟 Galois 扩张的情形：若 \(x \in L\) 使得存在 \(d \in A'\) 使 \(dx^{n} \in A'\)（对所有 \(n \geqslant 0\)），先证明可设 \(d \in A\)，然后推出：若 \(c_{i}\)（\(1 \leqslant i \leqslant m\)）是 x 在 K 上的极小多项式的系数，则 \(d^{m}c_{i}^{n} \in A\)（对所有 \(n \geqslant 0\)）。）

\begin{enumerate}
\def\labelenumi{\arabic{enumi}.}
\setcounter{enumi}{14}
\tightlist
\item
  设 K 为特征 0 的交换域，它包含所有单位根，并且存在 K 的 Galois 扩张，其 Galois 群不可解（例如参见《代数》第五章附录 I 第 1 节命题 2）。设 \(\Omega\) 为 K 的一个代数闭包，并设 \(\Omega_0\) 为满足下列条件的元素 \(z \in \Omega\) 组成的集：由 \(z\) 在 \(\Omega\) 中生成的 K 的 Galois 扩张具有可解的 Galois 群。
\end{enumerate}

\begin{enumerate}
\def\labelenumi{(\alph{enumi})}
\item
  证明 \(\Omega_0\) 是 \(\Omega\) 的一个异于 \(\Omega\) 的子群（参见《代数》第五章 § 10 第 4 节定理 1）。
\item
  设 A 为 \(\Omega[X]\) 的子环，由满足 \(\mathrm{P}(0) \in \Omega_{0}\) 的多项式 P 组成。证明 A 不是整闭的，但在其分式域 \(\Omega(X)\) 中，每个有某次幂属于 A 的元素本身属于 A。
\end{enumerate}

\begin{enumerate}
\def\labelenumi{\arabic{enumi}.}
\setcounter{enumi}{15}
\item
  设 B 为环，A 为 B 的子环，f 为 A 的理想，即 A-模 B/A 的零化子。设 U 为 \(\mathbf{X} = \operatorname{Spec}(\mathbf{A})\) 中闭集 \(\mathbf{V}(\mathfrak{f})\) 的余集，并设 \(u = {}^{a}\rho\)，其中 \(\rho: A \to B\) 是典范单射。证明 u 在 \({}^{-1}(U)\) 上的限制是 \({}^{-1}(U)\) 到 U 上的同胚（对每个元素 \(g \in f\)，考虑环 \(A_{g}\) 的谱，它与开集 \(X_{g} \subset U\) 等同）。
\item
  设 K 为域，B 为多项式环 \(K[X_{n}]_{n\in N}\)，A 为 B 的子环，由单项式 \(X_{n}^{2}\) 与 \(X_{n}^{3}\)（对所有 \(n\in N\)）生成；证明 B 是 A 的整闭包，且 B 在 A 中的导子化为 0；但若 \(S=A-\{0\}\)，则 \(S^{-1}B\) 在 \(S^{-1}A\) 中的导子不化为 0，且 \(S^{-1}A\) 是整闭的。
\item
  设 A 为整环，K 为其分式域，M 为有限生成 A-模，u 为 M 的一个自同态，\(u \otimes 1\) 为有限维 K-向量空间 \(M_{(K)} = M \otimes_{A} K\) 上相应的自同态。证明 \(u \otimes 1\) 的特征多项式的系数在 A 上整（利用第 1 节定理 1 的条件 \((E_{\mathrm{III}})\)）。
\item
  \begin{enumerate}
  \def\labelenumii{(\alph{enumii})}
  \tightlist
  \item
    设 A 为整闭整环，K 为其分式域，L 为 K 的可分代数扩张，x 为 L 中在 A 上整的一个元素，\(K'\) 为域 \(\mathrm{K}(x) \subset \mathrm{L}\)，F 为 x 在 K 上的极小多项式。若 x 在 K 上的次数为 n，证明 A 在 \(K'\) 中的整闭包包含在由元素 \(1/\mathrm{F}'(x)\)、\(x/\mathrm{F}'(x)\)、\(\ldots\)、\(x^{n-1}/\mathrm{F}'(x)\)（都属于 \(K'\)）生成的 A-模之中。（若 \(z \in K'\) 在 A 上整，写下
  \end{enumerate}
\end{enumerate}

\[
z \mathrm{F} ^ {\prime} (x) = g (x), \quad \text { where } \quad g (\mathbf {X}) = b _ {0} + b _ {1} \mathbf {X} + \dots + b _ {n - 1} \mathbf {X} ^ {n - 1} \in \mathbf {K} [ \mathbf {X} ].
\]

借助 Lagrange 插值公式证明

\[
g (\mathbf {X}) = \operatorname{Tr} _ {\mathbf {K} ^ {\prime} [ \mathbf {X} ] / \mathbf {K} [ \mathbf {X} ]} \left(z. \frac {\mathrm{F} (\mathbf {X})}{\mathbf {X} - x}\right) \cdot)
\]

\begin{enumerate}
\def\labelenumi{(\alph{enumi})}
\setcounter{enumi}{1}
\tightlist
\item
  设 A 为整闭整环，K 为其分式域，p 为 K 的特征指数。假设子环 \(A^{1/p}\)（\(K^{1/p}\) 中由 A 的元素的 p 次根组成的子环）是有限生成 A-模。证明对 K 的每个有限扩张 E，A 在 E 中的整闭包是有限生成 A-模（化归到 E 是拟 Galois 扩张的情形）。
\end{enumerate}

¶ * 20. 设 \(\mathbf{K}_0\) 为特征 \(p > 0\) 的完满域，\(\mathbf{K}\) 为有理函数域 \(\mathrm{K}_0(\mathrm{X}_n)_{n \in \mathbb{N}}\)，\(\mathbf{B}\) 为形式幂级数环 \(\mathrm{K}[[\mathrm{Z}]]\)，其中 \(\mathbf{Z}\) 是一个未定元；若 \(m\) 是 \(B\) 的极大理想，则 \(B\) 对 \(m\) -进拓扑是 Hausdorff 且完备的，并且是一个离散赋值环。设 \(L = K((Z))\) 为 \(B\) 的分式域，\(E_0\) 为 \(L\) 的子域，由 \(L^p\) 与元素 \(Z\) 和 \(X_n\)（对 \(n \in N\)）生成。

\begin{enumerate}
\def\labelenumi{(\alph{enumi})}
\item
  证明元素 \(c = \sum_{n=0}^{\infty} \mathbf{X}_n \mathbf{Z}^n\)（属于 \(L\)）不属于 \(E_0\)。
\item
  设 E 为 L 的包含 \(E_{0}\) 而不包含 c 的子域所成之集的一个极大元；记 \(A = B \cap E\)。证明 A 是一个离散赋值环（从而是 Noether 的），其分式域为 E；若 \(m'\) 是 A 的极大理想，则 \(m'^{k} = m^{k} \cap A\)，且 A 对 m-进拓扑在 B 中稠密；\(L = K(c)\) 是 K 的一个次数为 p 的扩张，B 是 A 在 L 中的整闭包，但 B 不是有限生成 A-模（利用第三章 §3 习题 9 (b)）。*
\end{enumerate}

\begin{enumerate}
\def\labelenumi{\arabic{enumi}.}
\setcounter{enumi}{20}
\tightlist
\item
  \begin{enumerate}
  \def\labelenumii{(\alph{enumii})}
  \tightlist
  \item
    设 \(K_{0}\) 为特征 p \textgreater{} 0 的完满域，L 为有理函数域 \(\mathrm{K}_{0}(\mathrm{X}_{n})_{n\in\mathbf{N}}\)，\(A'\) 为形式幂级数环 \(L[[Y,Z]]\)（其中 Y 与 Z 是两个未定元），A 为张量积环 \(K[[Y,Z]]\otimes_{K}L\)，这里记 \(K = L^{p} \subset L\)；A 是一个 Noether 局部环，它不完备，其完备化与 \(A'\) 等同（第三章 §3 习题 17），且 \(A'\) 在 A 上整。设 \(c_{n}\) 表示元素
  \end{enumerate}
\end{enumerate}

\[
\mathrm{X} _ {n} \mathrm{Y} + \mathrm{X} _ {n + 1} \mathrm{YZ} + \mathrm{X} _ {n + 2} \mathrm{YZ} ^ {2} + \dots
\]

（属于 A'）。设 B 为 A' 的由 A 与诸元素 \(c_{n}\)（\(n \geqslant 0\)）生成的子环。若 C = A{[}c₀{]}，则 C 是 Noether 的，且环 B（它不整闭）含于 C 的整闭包并包含 C；但证明 B 不是 Noether 的（注意 \(c_{n}\) 不属于 B 的由诸 \(c_{i}\)（其下标 i \textless{} n）生成的理想）。

\begin{enumerate}
\def\labelenumi{(\alph{enumi})}
\setcounter{enumi}{1}
\tightlist
\item
  设 \(p = 2\)；\(K_0\)、\(K\) 与 \(L\) 的意义与 (a) 中相同，这次考虑三个未定元的形式幂级数环 \(A' = L[[Y, Z, T]]\)，以及子环
\end{enumerate}

\[
\mathrm{A} = \mathrm{K} [ [ \mathrm{Y}, \mathrm{Z}, \mathrm{T} ] ] \otimes_ {\mathrm{K}} \mathrm{L},
\]

并写为

\[
b = \mathrm{Y} \sum_ {n \geqslant 0} \mathrm{X} _ {2 n} \mathrm{T} ^ {n} + \mathrm{Z} \sum_ {n \geqslant 0} \mathrm{X} _ {2 n + 1} \mathrm{T} ^ {n}.
\]

证明环 A{[}b{]} 是 Noether 的，但它的整闭包 B 不是 Noether 环。（由于 \(b^{2} \in A\)，A{[}b{]} 的分式域的每个元素形如 \(P + Qb\)，其中 P 与 Q 是诸 \(X_{k}\) 中有限个的线性组合，系数在 K{[}{[}Y, Z, T{]}{]} 的分式域中；为使这样一个元素在 A{[}b{]} 上整，证明其必要且充分条件是它属于 A'。证明 B 由诸元素

\[
b _ {i} = \mathrm{Y} \sum_ {n \geqslant 0} \mathrm{X} _ {2 (i + n)} \mathrm{T} ^ {n} + \mathrm{Z} \sum_ {n \geqslant 0} \mathrm{X} _ {2 (i + n) + 1} \mathrm{T} ^ {n} \quad \text { for } \quad i \geqslant 0
\]

生成，并像 (a) 那样完成论证。）

\begin{enumerate}
\def\labelenumi{\arabic{enumi}.}
\setcounter{enumi}{21}
\item
  设 A 为整闭整环，K 为其分式域；证明对 K 的每个扩张 L，存在 L 的一个整闭子整环 B，其分式域是 L 且满足 \(B \cap K = A\)（把 L 看作 K 的一个纯超越扩张 \(E = K(X_t)_{t \in I}\) 的代数扩张，从而利用习题 11 (a) 化归到 L 是 K 的代数扩张的情形）。
\item
  \begin{enumerate}
  \def\labelenumii{(\alph{enumii})}
  \tightlist
  \item
    设 K 为交换域，n 为与 K 的特征互素的整数，P 为 K{[}X{]} 中一个只有单根（在 K 的一个代数闭包中）的多项式。若记 \(f(\mathbf{X}, \mathbf{Y}) = \mathbf{Y}^{n} - \mathbf{P}(\mathbf{X})\)，证明：当 K 包含 n 次单位根时，整环 \(\mathbf{A} = \mathbf{K}[\mathbf{X}, \mathbf{Y}] / (f)\) 是整闭的，并且在其分式域中含有 K(X) 的由 f 的一个根 y 生成的可分扩张 L（f 视为 K(X){[}Y{]} 中的多项式）。（证明：若 L 的元素 z 在 K{[}X{]} 上整，则它属于 A；为此，把 z 写成形式
  \end{enumerate}
\end{enumerate}

\[
z = \sum_ {i = 0} ^ {n - 1} a _ {i} (\mathbf {X}) y ^ {i},
\]

其中 \(a_{i}(\mathbf{X}) \in \mathbf{K}(\mathbf{X})\)。证明诸元素 \(a_{i}(\mathbf{X})y^{i} (0 \leqslant i \leqslant n - 1)\) 在 \(\mathbf{K}[\mathbf{X}]\) 上整，并推出诸 \((a_{i}(\mathbf{X}))^{n}(\mathbf{P}(\mathbf{X}))^{i}\) 是 \(\mathbf{K}[\mathbf{X}]\) 的元素；从而诸 \(a_{i}(\mathbf{X})\) 也是。

\begin{enumerate}
\def\labelenumi{(\alph{enumi})}
\setcounter{enumi}{1}
\tightlist
\item
  设 K 为特征 \(p \neq 2\) 的非完满域；设 \(a \in K\) 为一个不属于 \(K^{p}\) 的元素，设 E 为域 \(\mathrm{K}(a^{1/p})\)，并记 \(\mathrm{P}(\mathrm{X}) = \mathrm{X}^{p} - a, f(\mathrm{X}, \mathrm{Y}) = \mathrm{Y}^{2} - \mathrm{P}(\mathrm{X})\)。若 A 为整闭整环 \(\mathrm{K}[\mathrm{X}, \mathrm{Y}]/(f)\)，证明 \(B = E \otimes_{K} A\) 是整环但不整闭（考虑 B 的分式域中的元素 \(y/(\mathrm{X} - a^{1/p})\)）。
\end{enumerate}

\begin{enumerate}
\def\labelenumi{\arabic{enumi}.}
\setcounter{enumi}{23}
\tightlist
\item
  设 \(\Delta\) 为一个以加法记的无扭交换群。设 A 为交换环，B 为交换 A-代数，\(h\colon\mathbf{A}\to\mathbf{B}\) 为定义 B 上代数结构的环同态。那么 \(h^{(\Delta)}\colon\mathbf{A}^{(\Delta)}\to\mathbf{B}^{(\Delta)}\)（《代数》
\end{enumerate}

第二章 § 11 习题 1）是代数 \(\mathbf{A}^{(\Delta)}\)（群 \(\Delta\) 关于 \(\mathbf{A}\) 的代数）到代数 \(\mathbf{B}^{(\Delta)}\)（群 \(\Delta\) 关于 \(\mathbf{B}\) 的代数）的一个（环）同态。设 \(\mathbf{P} = \sum_{\alpha \in \Delta} b_{\alpha} \otimes \mathbf{X}^{\alpha}\) 为 \(\mathbf{B}^{(\Delta)}\) 的一个元素。证明：\(\mathbf{P}\) 在 \(\mathbf{A}^{(\Delta)}\) 上整，当且仅当每个元素 \(b_{\alpha} \in \mathbf{B}\) 都在 \(\mathbf{A}\) 上整。（化归到 \(\Delta\) 是有限生成的情形，再化归到 \(\Delta = \mathbf{Z}\) 的情形。）证明：若 \(\mathbf{A}\) 是整环且整闭，则 \(\mathbf{A}^{(\Delta)}\) 也整闭。

\begin{enumerate}
\def\labelenumi{\arabic{enumi}.}
\setcounter{enumi}{24}
\tightlist
\item
  设 \(\Delta\) 为全序交换群。把第 8 节的命题 20 与 21 推广到 \(\Delta\) 型分次环（对命题 20 利用习题 24；对命题 21，把第 8 节的引理 4 推广到 \(\Delta\) 有限生成的情形，然后取正向极限证明命题 21）。
\end{enumerate}

¶ 26. (a) 设 A 为满足下列条件的整环：对 A 中每个 \(a \neq 0\)，理想 \(Aa\) 是一个族 \((q_t)\)（由 \(Aa\) 关于素理想 \(p_t\)（在包含 \(a\) 的素理想中极小者）的饱和组成）之交，且诸整环 \(A_{p_t}\) 是整闭的（相应地，完全整闭的）。证明 A 是整闭的（相应地，完全整闭的）。

\begin{enumerate}
\def\labelenumi{(\alph{enumi})}
\setcounter{enumi}{1}
\item
  设 A 为强 Lasker 局部环（第四章 § 2 习题 28），m 为其极大理想。假设对某个 \(a \in \mathbf{A}\)，m 是与 A/Aa 弱伴随的一个浸入素理想（第四章 § 1 习题 17）。证明 Aa: m 含于每个对与 A/Aa 弱伴随的素理想 p ≠ m 都准素的理想之中（对这样一个准素理想 q，考虑传递子 q: m，参见第四章 § 2 习题 30）。若进一步设 A 是整环，证明关系式 p.~(Aa: m) = Aa（对在包含 Aa 的素理想中极小的素理想 p）不可能成立（若 q 是 A 关于 p 的饱和，注意上述关系将蕴含在 A \(_{p}\) 中 qA \(_{p}\) = qpA \(_{p}\)，并借助第四章 § 2 习题 29 (d) 完成证明）。
\item
  证明：若 A 是强 Lasker 整环，且存在 A 的元素 \(a \neq 0\) 使得有与 A/Aa 弱伴随并且浸入的素理想，则 A 不是完全整闭的。（化归到 (b) 中考虑的情形；用 (b) 的记号，那时存在 \(b \in \mathbf{A}a\) : m 使得 \(b \notin \mathbf{A}a\) 且 \(bp \subset am\)；对 \(n\) 用归纳法推出 \(b^n p \subset a^n m\)（对所有 \(n\)）。）
\item
  特别地，由 (a) 与 (c) 推出：Noether 整环是整闭的，当且仅当对 A 中所有 \(a \neq 0\)，与 A/Aa 伴随的素理想 \(\mathfrak{p}\) 都不是浸入的，且 \(\mathrm{A}_{\mathfrak{p}}\) 是整闭的（参见第七章 § 1 第 4 节命题 8）。
\end{enumerate}

\begin{enumerate}
\def\labelenumi{\arabic{enumi}.}
\setcounter{enumi}{26}
\tightlist
\item
  设 A 为整闭但非完全整闭的整环，K 为其分式域；设 \(a \neq 0\) 为 A 的一个不可逆元素，对它存在 A 中的 \(d \neq 0\) 使 \(da^{-n} \in A\)（对每个整数 \(n \geqslant 0\)）。
\end{enumerate}

证明在形式幂级数域 \(\mathbf{K}((\mathbf{X}))\) 中存在一个形式幂级数 \(\sum_{k=1}^{\infty} u_k \mathbf{X}^k = f(\mathbf{X})\)，满足方程

\[
(f (\mathbf {X})) ^ {2} - a f (\mathbf {X}) + \mathbf {X} = 0,
\]

从而它在 A{[}{[}X{]}{]} 上整，属于 A{[}{[}X{]}{]} 的分式域但不属于 A{[}{[}X{]}{]}，故 A{[}{[}X{]}{]} 不是整闭的。

\begin{enumerate}
\def\labelenumi{\arabic{enumi}.}
\setcounter{enumi}{27}
\item
  设 \(k\) 为域，\(\mathbf{A}_0\) 为两个无穷未定元族的多项式环 \(k[\mathrm{X}_n, \mathrm{Y}_n]_{n \in \mathbf{Z}}\)。设 \(G\) 为一个无穷单生成群（从而同构于 \(\mathbf{Z}\)），并设 \(\sigma\) 为 \(\mathcal{G}\) 的一个生成元。\(G\) 定义为在 \(\mathbf{A}_0\) 上作用，其条件为 \(\sigma(\mathrm{Y}_n) = \mathrm{Y}_n\) 与 \(\sigma(\mathrm{X}_n) = \mathrm{X}_{n+1}\)（对所有 \(n \in \mathbf{Z}\)）。设 \(\Im\) 为 \(\mathbf{A}_0\) 中由元素 \(\mathrm{Y}_n(\mathrm{X}_n - \mathrm{X}_{n+1})\)（对所有 \(n \in \mathbf{Z}\)）生成的理想，并设 \(A\) 为商环 \(A_0 / \Im\)；由于 \(\sigma(\Im) \subset \Im\)，群 \(\mathcal{G}\) 通过过渡到商也在 \(A\) 上作用。设 \(S\) 为 \(A\) 中由诸 \(Y_n\) 在 \(A\) 中的典范像生成的乘性子集，于是 \(S\) 由在 \(G\) 下不变的元素组成。证明 \(S^{-1}(A^{\mathcal{G}}) \neq (S^{-1}A)^{\mathcal{G}}\)。
\item
  设 k 为特征 \(\neq2\) 的域，A 为一个未定元的多项式环 k{[}T{]}，a、b 为 k 中两个元素，满足 \(a\neq0\)、\(b\neq0\) 且 \(a\neq b\)。设 \(\mathbf{K}=k(\mathbf{T})\) 为 A 的分式域，L、M 为 K 的扩张，分别由 \(\mathbf{X}^{2}-\mathbf{T}(\mathbf{T}-a)\) 的一个根和 \(\mathbf{X}^{2}-\mathbf{T}(\mathbf{T}-b)\) 的一个根添加到 K 而得。设 B、C 分别为 A 在 L 与 M 中的整闭包，它们是有限生成自由 A-模。证明 \(L\otimes_{K}M\) 是一个域，是 K 的可分有限扩张，但 \(B\otimes_{A}C\)（它与 \(L\otimes_{K}M\) 的一个子环典范地等同）不是 A 在 \(L\otimes_{K}M\) 中的整闭包。
\end{enumerate}
